% ===========================================================================
%  Documento principal — Tesis de Diploma
%  Título: Un lenguaje de dominio específico para la generación declarativa
%          de hojas de cálculo de planificación
%  Autor:  Jose Morales Lazo
%  Tutor:  MSc. Fernando Rodríguez Flores
%  MATCOM — Universidad de La Habana, 2026
% ===========================================================================
\documentclass[12pt,a4paper]{report}

% ---- Codificación e idioma -------------------------------------------------
\usepackage[utf8]{inputenc}
\usepackage[T1]{fontenc}
\usepackage[spanish,es-noquoting,es-noshorthands]{babel}

% ---- Geometría -------------------------------------------------------------
\usepackage{geometry}
\geometry{left=3cm,right=2.5cm,top=3cm,bottom=2.5cm}

% ---- Tipografía ------------------------------------------------------------
\usepackage{lmodern}
\usepackage{microtype}
\usepackage{parskip}

% ---- Matemáticas -----------------------------------------------------------
\usepackage{amsmath}
\usepackage{amssymb}

% ---- Tablas ----------------------------------------------------------------
\usepackage{booktabs}
\usepackage{array}
\usepackage{longtable}
\usepackage{multirow}

% ---- Gráficos --------------------------------------------------------------
\usepackage{graphicx}
\usepackage{float}
\usepackage{caption}

% ---- Colores ---------------------------------------------------------------
\usepackage{xcolor}
\definecolor{lispkw}{rgb}{0.0,0.35,0.65}
\definecolor{lispstr}{rgb}{0.13,0.55,0.13}
\definecolor{lispcomment}{rgb}{0.5,0.5,0.5}
\definecolor{lispbg}{rgb}{0.97,0.97,0.99}
\definecolor{pybg}{rgb}{0.97,0.99,0.97}
\definecolor{nodeborder}{rgb}{0.3,0.45,0.7}
\definecolor{sectioncolor}{rgb}{0.15,0.25,0.5}
\definecolor{graylight}{rgb}{0.94,0.94,0.94}

% ---- Listados de código ----------------------------------------------------
\usepackage{listings}

\lstdefinelanguage{CommonLisp}{
  keywords={defclass,defmacro,defun,defparameter,defvar,defgeneric,defmethod,
            progn,let,let*,loop,for,collect,append,list,when,unless,if,cond,
            load,format,lambda,provide,values,ecase,mapcar,
            def-table,conditional-rendering,exists,col,tcol,trange,
            countif,counta,sum-range,str,nav,ultima-fila,primera-fila,
            libro,hoja,hoja-v,tabla,xl-generate,xl-run-generated,
            add,subtract,multiply,divide,equals,different,
            non-empty,previous-row,next-of,previous-of,show-nothing,
            concat,time-add,range,lookup,sheet-ref,cross-cell,
            source-row,turno-aula-row,sheet-id,collect-over,param,
            defclass*,symb,mkstr,keywd},
  keywordstyle=\color{lispkw}\bfseries,
  sensitive=false,
  comment=[l]{;;},
  commentstyle=\color{lispcomment}\itshape,
  morestring=[b]",
  stringstyle=\color{lispstr},
  basicstyle=\ttfamily\small,
  backgroundcolor=\color{lispbg},
  frame=single, framerule=0.4pt,
  rulecolor=\color{nodeborder!60},
  breaklines=true, showstringspaces=false,
  tabsize=2, xleftmargin=1em, xrightmargin=0.5em,
  aboveskip=0.8em, belowskip=0.5em,
}

\lstdefinestyle{pythonstyle}{
  language=Python,
  basicstyle=\ttfamily\small,
  backgroundcolor=\color{pybg},
  keywordstyle=\color{lispkw}\bfseries,
  stringstyle=\color{lispstr},
  commentstyle=\color{lispcomment}\itshape,
  frame=single, framerule=0.4pt,
  rulecolor=\color{nodeborder!60},
  breaklines=true, showstringspaces=false,
  tabsize=4, xleftmargin=1em, xrightmargin=0.5em,
  aboveskip=0.8em, belowskip=0.5em,
}

\lstset{language=CommonLisp}

% ---- Formato de secciones --------------------------------------------------
\usepackage{titlesec}
\titleformat{\chapter}[display]
  {\normalfont\Large\bfseries\color{sectioncolor}}
  {\chaptertitlename\ \thechapter}{16pt}{\Huge}
\titleformat{\section}
  {\normalfont\large\bfseries\color{sectioncolor}}{\thesection}{1em}{}
\titleformat{\subsection}
  {\normalfont\normalsize\bfseries\color{sectioncolor!80}}{\thesubsection}{1em}{}

% ---- Enumeraciones ---------------------------------------------------------
\usepackage{enumitem}

% ---- Teoremas y definiciones -----------------------------------------------
\usepackage{amsthm}
\newtheoremstyle{defstyle}{6pt}{6pt}{\itshape}{}{\bfseries}{.}{0.5em}{}
\theoremstyle{defstyle}
\newtheorem{definition}{Definici\'on}[chapter]

% ---- Hipervínculos ---------------------------------------------------------
\usepackage[hidelinks]{hyperref}
\hypersetup{
  pdftitle={Un DSL para la generaci\'on declarativa de hojas de c\'alculo de planificaci\'on},
  pdfauthor={Jose Morales Lazo},
}

% ---- Encabezados -----------------------------------------------------------
\usepackage{fancyhdr}
\pagestyle{fancy}
\fancyhf{}
\fancyhead[LE,RO]{\thepage}
\fancyhead[RE]{\nouppercase{\leftmark}}
\fancyhead[LO]{\nouppercase{\rightmark}}
\renewcommand{\headrulewidth}{0.4pt}

% ---- Utilidades ------------------------------------------------------------
\newcommand{\dsl}[1]{\texttt{#1}}
\newcommand{\file}[1]{\texttt{#1}}
\newcommand{\xlsx}{\texttt{.xlsx}}

% ===========================================================================
\begin{document}

% ---------------------------------------------------------------------------
%  PORTADA
% ---------------------------------------------------------------------------
\begin{titlepage}
\centering
{\large Universidad de La Habana\\
Facultad de Matem\'atica y Computaci\'on}

\vspace{2cm}
% \includegraphics[width=3cm]{uh-logo}  % añadir logo si disponible
\vspace{2cm}

{\LARGE\bfseries Un lenguaje de dominio espec\'{\i}fico para la\\[6pt]
generaci\'on declarativa de hojas de c\'alculo\\[6pt]
de planificaci\'on}

\vspace{2.5cm}

{\large Autor: \textbf{Jose Morales Lazo}}

\vspace{0.8cm}

{\large Tutor:\\
\textbf{MSc. Fernando Rodr\'{\i}guez Flores}}

\vspace{2cm}

{\normalsize Trabajo de Diploma\\
presentado en opci\'on al t\'{\i}tulo de\\
Licenciado en Ciencia de la Computaci\'on}

\vfill
{\normalsize La Habana, junio 2026}
\end{titlepage}

% ---------------------------------------------------------------------------
%  FRASE
% ---------------------------------------------------------------------------
\thispagestyle{empty}
\vspace*{8cm}
\begin{flushright}
\itshape ``La simplicidad es el prerrequisito de la fiabilidad.''\\[4pt]
\upshape Edsger W. Dijkstra
\end{flushright}
\newpage

% ---------------------------------------------------------------------------
%  AGRADECIMIENTOS
% ---------------------------------------------------------------------------
\chapter*{Agradecimientos}
\addcontentsline{toc}{chapter}{Agradecimientos}
\thispagestyle{plain}

A mi tutor, el MSc. Fernando Rodr\'{\i}guez Flores, por su paciencia, sus ideas
y su entusiasmo permanente ante los lenguajes y las hojas de c\'alculo.

A mi familia, por el apoyo incondicional a lo largo de la carrera.

A los colegas y profesores de MATCOM que convirtieron cada problema
acad\'emico en una oportunidad de automatizarlo.

\newpage

% ---------------------------------------------------------------------------
%  OPINIÓN DEL TUTOR
% ---------------------------------------------------------------------------
\chapter*{Opini\'on del Tutor}
\addcontentsline{toc}{chapter}{Opini\'on del Tutor}
\thispagestyle{plain}

La gesti\'on de la informaci\'on en entornos acad\'emicos genera, de forma
recurrente, documentos de planificaci\'on tabular: horarios de docencia,
calendarios de defensas, distribuciones de aulas. En todos estos casos, la
herramienta elegida es la hoja de c\'alculo. Nada de esto sorprende; lo que
s\'{\i} resulta costoso es que cada vez que se quiere automatizar la creaci\'on
de esa hoja se repite el mismo patr\'on: c\'odigo imperativo que mezcla la
l\'ogica del problema con los detalles de la \textit{API} de la biblioteca
de salida.

El trabajo que se presenta en este documento aborda precisamente ese costo.
Jose Morales Lazo dise\~n\'o e implement\'o un lenguaje de dominio espec\'{\i}fico
(DSL) interno en Common Lisp que permite declarar tablas, columnas calculadas,
reglas de visualizaci\'on condicional y relaciones entre tablas, sin escribir
una sola letra de columna Excel, un solo \'{\i}ndice de fila ni una cadena de
f\'ormula. El compilador del DSL, implementado como un backend de funciones
gen\'ericas CLOS, produce autom\'aticamente el programa Python que construye
el archivo \xlsx{} con todas las f\'ormulas y reglas de formato condicional
incrustadas.

El diplomante demostr\'o durante el trabajo una notable capacidad para
identificar la abstracci\'on correcta, dise\~narla y construirla con rigor.
Los tres casos de uso reales validados ---horario de defensas de tesis,
horario docente de la facultad y horario de programaci\'on televisiva---
muestran que el lenguaje es suficientemente expresivo para cubrir problemas
de complejidad creciente.

Estoy muy satisfecho con los resultados obtenidos y considero que el trabajo
re\'une los m\'eritos suficientes para optar al t\'{\i}tulo de Licenciado en
Ciencia de la Computaci\'on.

\vspace{2cm}
\begin{flushright}
\rule{6cm}{0.4pt}\\
MSc. Fernando Rodr\'{\i}guez Flores\\
Facultad de Matem\'atica y Computaci\'on\\
Universidad de La Habana\\
Junio, 2026
\end{flushright}
\newpage

% ---------------------------------------------------------------------------
%  RESUMEN
% ---------------------------------------------------------------------------
\chapter*{Resumen}
\addcontentsline{toc}{chapter}{Resumen}
\thispagestyle{plain}

La creaci\'on autom\'atica de hojas de c\'alculo de planificaci\'on exige, en
la pr\'actica, c\'odigo imperativo que entremezcla la l\'ogica del problema
---qui\'en tiene conflicto de horario, cu\'antas defensas acumula un profesor---
con los detalles mec\'anicos de la API de salida: letras de columna Excel,
\'{\i}ndices de fila absolutos, cadenas de f\'ormula, llamadas a la biblioteca
de formato condicional.  Esta mezcla eleva el costo de desarrollo y dificulta
la reutilizaci\'on en problemas similares.

El objetivo de este trabajo es separar ambas capas mediante un lenguaje de
dominio espec\'{\i}fico (DSL) interno en Common~Lisp.  El programador declara
\emph{qu\'e} debe contener el libro ---tablas, columnas calculadas, cuantificadores
existenciales, reglas de coloreado, totales relativos--- y un compilador basado en
el Sistema de Objetos de Common Lisp (CLOS) traduce esa declaraci\'on a un
programa Python que construye el archivo \xlsx{}.  Las reglas de formato
condicional se incrustan como \texttt{FormulaRule} y se reevaluan autom\'aticamente
en Excel al editar los datos, sin necesidad de volver a ejecutar el DSL.

La validaci\'on sobre tres casos reales de la Facultad de Matem\'atica y
Computaci\'on de la Universidad de La Habana confirma que el lenguaje captura
los patrones recurrentes de la planificaci\'on tabular con una cantidad
de c\'odigo significativamente menor que la alternativa imperativa.

\medskip\noindent
\textbf{Palabras clave:} Lenguaje de dominio espec\'{\i}fico; Hoja de c\'alculo;
Generaci\'on autom\'atica de c\'odigo; Common Lisp; CLOS; Macros;
Formato condicional; openpyxl.
\newpage

% ---------------------------------------------------------------------------
%  ABSTRACT
% ---------------------------------------------------------------------------
\chapter*{Abstract}
\addcontentsline{toc}{chapter}{Abstract}
\thispagestyle{plain}

Automatically generating planning spreadsheets requires, in practice,
imperative code that interleaves domain logic---who has a scheduling
conflict, how many thesis defenses a professor is involved in---with
the mechanical details of the output API: Excel column letters, absolute
row indices, formula strings, and conditional formatting library calls.
This interleaving raises development costs and hampers reuse across
similar problems.

The objective of this work is to separate these two layers through an
internal domain-specific language (DSL) built in Common~Lisp.  The
programmer declares \emph{what} the workbook should contain---tables,
computed columns, existential quantifiers, coloring rules, relative
totals---and a compiler based on the Common Lisp Object System (CLOS)
translates that declaration into a Python program that builds the
\xlsx{} file.  Conditional formatting rules are embedded as
\texttt{FormulaRule} objects and are automatically re-evaluated in
Excel when data is edited, without re-running the DSL.

Validation on three real use cases from the Faculty of Mathematics and
Computer Science of the University of Havana confirms that the language
captures the recurring patterns of tabular planning with significantly
less code than the imperative alternative.

\medskip\noindent
\textbf{Keywords:} Domain-specific language; Spreadsheet;
Automatic code generation; Common Lisp; CLOS; Macros;
Conditional formatting; openpyxl.
\newpage

% ---------------------------------------------------------------------------
%  ÍNDICE GENERAL
% ---------------------------------------------------------------------------
\tableofcontents
\newpage

% ---------------------------------------------------------------------------
%  CUERPO PRINCIPAL
% ---------------------------------------------------------------------------
% ===========================================================================
%  Introducción
% ===========================================================================
\chapter*{Introducci\'on}
\addcontentsline{toc}{chapter}{Introducci\'on}
\markboth{Introducci\'on}{Introducci\'on}
\label{cap:intro}

% ---------------------------------------------------------------------------
\section*{Contexto}
% ---------------------------------------------------------------------------

La gesti\'on de la informaci\'on acad\'emica en una facultad universitaria
produce, de manera recurrente, un tipo particular de documento: la hoja de
c\'alculo de planificaci\'on.  Los horarios de docencia, los calendarios de
defensas de tesis, los cuadros de distribuci\'on de aulas y las
programaciones de actividades extracurriculares comparten una estructura
com\'un: tablas relacionadas, con columnas que resumen informaci\'on de
otras tablas mediante f\'ormulas, y con celdas que se colorean
autom\'aticamente cuando se detecta un conflicto o una condici\'on
relevante~\cite{ldmag2025}.

La hoja de c\'alculo es la herramienta elegida para estos documentos por
razones pr\'acticas: est\'a disponible en cualquier puesto de trabajo, no
requiere instalaci\'on de software adicional y produce un resultado que el
usuario final puede editar directamente.  Sin embargo, cuando el documento
debe generarse autom\'aticamente a partir de datos externos ---una base de
datos, un formulario, un archivo JSON--- la simplicidad operativa de la hoja
de c\'alculo se convierte en complejidad de programaci\'on.

% ---------------------------------------------------------------------------
\section*{El problema}
% ---------------------------------------------------------------------------

El problema central que este trabajo aborda es la mezcla de responsabilidades
en el c\'odigo que genera hojas de c\'alculo complejas.

Para ilustrarlo, consideremos un caso concreto: construir autom\'aticamente
un horario de defensas de tesis en Excel.  El documento debe incluir una
tabla de defensas con las columnas de rol del tribunal, una tabla de
profesores con el n\'umero de defensas en que participa cada uno, y una regla
de formato condicional que coloree en rojo a los profesores que aparecen en
dos defensas programadas en el mismo slot horario.

Un programa imperativo que genera este documento en Python usando la
biblioteca \texttt{openpyxl} debe conocer simult\'aneamente:
\begin{itemize}
  \item que la columna \texttt{tutor} ocupa la letra B y \texttt{secretario}
    ocupa la letra F en la hoja de c\'alculo concreta;
  \item que los datos de defensas empiezan en la fila 4 y terminan en la fila 6;
  \item que la f\'ormula de conteo es
    \texttt{COUNTIF(\$B\$4:\$F\$6,K4)};
  \item que la regla de coloreado condicional es una \texttt{FormulaRule}
    con la expresi\'on
    \texttt{SUMPRODUCT((\$B\$4:\$B\$6=\$K4)+\ldots) * COUNTIFS(\ldots) > 0 > 0}.
\end{itemize}

Si el n\'umero de defensas cambia, o si se agrega una columna de rol, o si
las tablas se reorganizan, \emph{todos} esos valores deben recalcularse
manualmente.  El c\'odigo no expresa la intenci\'on ---``colorear a los
profesores con conflicto de horario''--- sino el mecanismo concreto de su
implementaci\'on en una versi\'on espec\'{\i}fica de la hoja.

Esta situaci\'on no es exclusiva del horario de defensas.  El horario docente
de la facultad, que distribuye grupos, turnos y aulas en m\'ultiples hojas
con referencias cruzadas, y el horario de programaci\'on televisiva, que
calcula el fin de cada programa a partir de su hora de inicio y duraci\'on,
presentan el mismo patr\'on: la l\'ogica del problema queda ahogada entre
detalles de la API.

La pregunta de investigaci\'on central de este trabajo es:

\medskip
\begin{center}
\itshape
?`Puede un lenguaje de dominio espec\'{\i}fico separar la intenci\'on
sem\'antica del mecanismo de implementaci\'on, permitiendo generar
hojas de c\'alculo complejas de planificaci\'on con un c\'odigo declarativo
de nivel de abstracci\'on superior?
\end{center}
\medskip

% ---------------------------------------------------------------------------
\section*{Hip\'otesis de la investigaci\'on}
% ---------------------------------------------------------------------------

La utilizaci\'on de un lenguaje de dominio espec\'{\i}fico (DSL) interno,
implementado mediante macros de Common Lisp y un compilador basado en el
Sistema de Objetos de Common Lisp (CLOS), permite separar la descripci\'on
sem\'antica de una hoja de c\'alculo de planificaci\'on ---su estructura
tabular, sus relaciones, sus f\'ormulas y sus condiciones de coloreado--- de
los detalles concretos de la API de generaci\'on de archivos \xlsx{}.  Esta
separaci\'on reduce el c\'odigo necesario, facilita la reutilizaci\'on ante
cambios en los datos y hace que el programa resultante sea legible por el
experto del dominio sin conocimiento de la API de destino.

% ---------------------------------------------------------------------------
\section*{Objetivo general}
% ---------------------------------------------------------------------------

Dise\~nar e implementar un lenguaje de dominio espec\'{\i}fico en Common Lisp
para la declaraci\'on y generaci\'on autom\'atica de hojas de c\'alculo de
planificaci\'on, incluyendo tablas relacionadas, columnas calculadas,
cuantificaci\'on existencial y formato condicional persistente en el archivo
\xlsx{}.

% ---------------------------------------------------------------------------
\section*{Objetivos espec\'{\i}ficos}
% ---------------------------------------------------------------------------

\begin{enumerate}
  \item Analizar las caracter\'{\i}sticas recurrentes de las hojas de c\'alculo
    de planificaci\'on en el contexto acad\'emico, identificando los patrones
    comunes de estructura, c\'alculo y visualizaci\'on condicional.

  \item Dise\~nar la sintaxis y la sem\'antica de un DSL declarativo capaz de
    expresar tablas relacionadas, referencias cruzadas entre tablas,
    columnas computadas, cuantificaci\'on existencial y posicionamiento
    relativo de expresiones auxiliares.

  \item Definir un \'arbol de sintaxis abstracta (AST) cuyas clases
    representen los conceptos del dominio de manera independiente del
    backend de generaci\'on.

  \item Implementar el DSL como un conjunto de macros de Common Lisp que
    construyen nodos del AST en tiempo de expansi\'on de macros.

  \item Implementar un backend de generaci\'on de c\'odigo basado en una
    funci\'on gen\'erica CLOS que traduzca cada tipo de nodo AST a
    instrucciones Python para la biblioteca openpyxl.

  \item Validar el lenguaje sobre tres casos de uso reales de la Facultad
    de Matem\'atica y Computaci\'on de la Universidad de La Habana: el
    horario de defensas de tesis, el horario docente de la facultad y
    el horario de programaci\'on televisiva.
\end{enumerate}

% ---------------------------------------------------------------------------
\section*{Estructura del documento}
% ---------------------------------------------------------------------------

El documento se organiza de la siguiente manera.

El cap\'{\i}tulo~\ref{cap:teoria} sienta las bases te\'oricas y conceptuales del
trabajo.  Se presentan las definiciones fundamentales sobre hojas de c\'alculo,
se introduce el concepto de lenguaje de dominio espec\'{\i}fico y su uso en la
construcci\'on de compiladores, se explica el papel de los \'arboles de sintaxis
abstracta como representaci\'on intermedia, y se describen las caracter\'{\i}sticas
de Common Lisp ---macros, CLOS y despacho m\'ultiple--- que hacen de \'el una
plataforma id\'onea para el dise\~no del DSL.

El cap\'{\i}tulo~\ref{cap:dominio} analiza el dominio de la planificaci\'on
tabular a partir de tres casos reales.  Se identifican los elementos comunes
en estos documentos y se formulan los requisitos concretos que el lenguaje debe
satisfacer.

El cap\'{\i}tulo~\ref{cap:dsl} describe el lenguaje de dominio espec\'{\i}fico
dise\~nado.  Se presentan el \'arbol de sintaxis abstracta con sus cuatro
familias de nodos, las formas del DSL visible al programador y su sem\'antica,
y el flujo completo desde la declaraci\'on hasta la producci\'on del
archivo \xlsx{}.

El cap\'{\i}tulo~\ref{cap:generacion} detalla el proceso de generaci\'on de
c\'odigo.  Se explica c\'omo el backend resuelve las coordenadas Excel, c\'omo
compila cada tipo de expresi\'on a cadenas de f\'ormula, y c\'omo genera las
f\'ormulas SUMPRODUCT que implementan el formato condicional persistente.

El cap\'{\i}tulo~\ref{cap:implementacion} describe los detalles de implementaci\'on:
la organizaci\'on en cuatro archivos con responsabilidades separadas, la macro
\dsl{defclass*} que automatiza la definici\'on de nodos del AST, y los
principios de extensibilidad del sistema.

El documento concluye con las conclusiones del trabajo, las recomendaciones
para investigaciones futuras y las referencias bibliogr\'aficas consultadas.

% ===========================================================================
%  Capítulo 1 — Preliminares teóricos
% ===========================================================================
\chapter{Preliminares te\'oricos}
\label{cap:teoria}

En este cap\'{\i}tulo se presentan los conceptos y herramientas que fundamentan
el trabajo.  La secci\'on~\ref{sec:hoja-calculo} define los t\'erminos clave
del dominio de las hojas de c\'alculo.  La secci\'on~\ref{sec:planificacion}
sit\'ua ese dominio en el contexto acad\'emico.  Las secciones~\ref{sec:dsl}
y~\ref{sec:ast} introducen los lenguajes de dominio espec\'{\i}fico y los
\'arboles de sintaxis abstracta.  Finalmente, la secci\'on~\ref{sec:lisp}
describe las caracter\'{\i}sticas de Common Lisp que lo convierten en la
plataforma de implementaci\'on elegida.

% ---------------------------------------------------------------------------
\section{Hojas de c\'alculo}
\label{sec:hoja-calculo}
% ---------------------------------------------------------------------------

Una \textbf{hoja de c\'alculo} (\textit{spreadsheet}) es una aplicaci\'on
para organizar, calcular y visualizar datos en una cuadr\'{\i}cula
bidimensional de celdas.  Cada celda puede contener un valor literal
---un n\'umero, una fecha o una cadena de texto--- o una
\textbf{f\'ormula}: una expresi\'on cuyo resultado se recalcula
autom\'aticamente cuando cambia cualquier celda de la que depende~\cite{microsoft-excel}.

\begin{description}
  \item[Workbook (libro de trabajo).] Un archivo \xlsx{} que agrupa una
    o m\'as hojas de trabajo.  Es la unidad de intercambio del formato
    Office Open XML~\cite{ooxml}.

  \item[Hoja de trabajo.] Una cuadr\'{\i}cula de celdas dentro del workbook,
    identificada por un nombre (p.~ej., \texttt{Defensas}).

  \item[Celda.] La unidad m\'{\i}nima de almacenamiento, referenciada por
    una direcci\'on del tipo \texttt{B4} (columna B, fila 4) o
    \texttt{\$B\$4} (referencia absoluta).

  \item[F\'ormula.] Una expresi\'on que comienza con \texttt{=} y puede
    contener referencias a otras celdas, rangos y funciones predefinidas
    como \texttt{COUNTIF}, \texttt{COUNTA}, \texttt{SUM}, \texttt{IF} o
    \texttt{SUMPRODUCT}.

  \item[Rango.] Un bloque rectangular de celdas contiguas, expresado como
    \texttt{\$B\$4:\$F\$6} (columnas B--F, filas 4--6).

  \item[Formato condicional.] Un mecanismo que aplica estilos visuales
    ---color de relleno, color de texto, borde--- a las celdas de un rango
    cuando se cumple una condici\'on evaluada en tiempo de apertura o
    edici\'on del archivo.  Una \textbf{FormulaRule} es el tipo m\'as general
    de regla de formato condicional: la condici\'on es una f\'ormula
    arbitraria, lo que permite expresar condiciones que involucran otras
    hojas, otras celdas o funciones de agregaci\'on~\cite{openpyxl}.
\end{description}

La biblioteca \textbf{openpyxl} es la herramienta Python elegida para
producir archivos \xlsx{} en este trabajo.  Permite crear workbooks,
escribir valores y f\'ormulas en celdas, y definir \texttt{FormulaRule}
que quedan incrustadas en el archivo y se eval\'uan en Excel o LibreOffice
Calc sin necesidad de c\'odigo adicional~\cite{openpyxl}.

% ---------------------------------------------------------------------------
\section{Planificaci\'on tabular en el \'ambito acad\'emico}
\label{sec:planificacion}
% ---------------------------------------------------------------------------

La planificaci\'on tabular agrupa los problemas de organizaci\'on cuya
soluci\'on se representa de forma natural como una o m\'as tablas
relacionadas.  En el contexto acad\'emico, los ejemplos m\'as frecuentes son:

\begin{itemize}
  \item \textbf{Horarios de docencia.} Distribuci\'on de grupos, asignaturas,
    profesores y aulas en slots de d\'{\i}a y hora a lo largo de la semana.
    Cada columna calcula la ocupaci\'on de un aula a partir de las asignaciones
    de los distintos grupos.

  \item \textbf{Planificaci\'on de defensas de tesis.} Asignaci\'on de
    estudiantes, fechas, horarios y miembros de tribunal (tutor, oponente,
    presidente, vocal, secretario) a defensas individuales.  El requisito
    cr\'{\i}tico es detectar conflictos: un mismo profesor en dos defensas
    del mismo slot horario.

  \item \textbf{Asignaci\'on de docencia.} Relaci\'on entre profesores,
    asignaturas y grupos, con conteos de horas acumuladas y alertas cuando
    se supera un umbral.
\end{itemize}

Estos documentos comparten una estructura t\'{\i}pica~\cite{ldmag2025}:
\begin{enumerate}
  \item Dos o m\'as tablas relacionadas que comparten columnas de clave
    (p.~ej., nombre de profesor, slot horario).
  \item Columnas calculadas en una tabla que agregan datos de otra
    (conteos, sumas, b\'usquedas).
  \item Reglas de formato condicional que colorean filas cuando se detecta
    una condici\'on problem\'atica (conflicto de horario, desequilibrio
    de carga).
  \item Totales y etiquetas fuera de las tablas, cuya posici\'on depende del
    n\'umero de filas de datos, que cambia en cada ejecuci\'on.
\end{enumerate}

Estos cuatro patrones son los que el DSL debe capturar de forma declarativa.

% ---------------------------------------------------------------------------
\section{Lenguajes de dominio espec\'{\i}fico}
\label{sec:dsl}
% ---------------------------------------------------------------------------

Un \textbf{lenguaje de dominio espec\'{\i}fico} (\textit{Domain-Specific
Language}, DSL) es un lenguaje de programaci\'on o especificaci\'on dise\~nado
para expresar soluciones dentro de un dominio de problemas concreto, en
contraposici\'on a los lenguajes de prop\'osito general que aspiran a cubrir
cualquier problema~\cite{fowler2010}.

\begin{definition}[DSL, \cite{van-deursen2000}]
Un lenguaje de dominio espec\'{\i}fico es un lenguaje de programaci\'on o
especificaci\'on ejecutable que ofrece, mediante abstracciones apropiadas,
poder expresivo centrado en y normalmente restringido a un dominio de
problemas particular.
\end{definition}

Los DSL se clasifican en dos grandes tipos seg\'un su relaci\'on con el
lenguaje anfitri\'on~\cite{fowler2010}:

\begin{description}
  \item[DSL externo.] Tiene su propia sintaxis independiente, requiere un
    lexer y un parser propios, y suele estar embebido como datos en otro
    programa (p.~ej., SQL, HTML, \LaTeX).

  \item[DSL interno (embebido).] Se expresa directamente en el lenguaje
    anfitri\'on mediante una API o un sistema de macros que produce c\'odigo
    v\'alido del anfitri\'on~\cite{mernik2005}.  El programador usa el
    lenguaje anfitri\'on como plataforma; el DSL es una biblioteca con
    vocabulario de dominio.
\end{description}

El DSL desarrollado en este trabajo es \textbf{interno}: est\'a escrito en
Common Lisp, sus formas son expresiones Lisp v\'alidas y no requiere ning\'un
parser externo.  Los macros de Lisp operan como el compilador del DSL,
expandiendo las formas del dominio en constructores de nodos del AST.

La ventaja principal de los DSL internos en Lisp es que el lenguaje anfitri\'on
es homoic\'onico: el c\'odigo tiene la misma estructura que los datos
(listas anidadas), lo que hace que escribir transformaciones de c\'odigo
---es decir, macros--- sea tan natural como escribir funciones ordinarias~\cite{graham1993}.

% ---------------------------------------------------------------------------
\section{\'Arboles de sintaxis abstracta}
\label{sec:ast}
% ---------------------------------------------------------------------------

Un \textbf{\'arbol de sintaxis abstracta} (AST, del ingl\'es
\textit{Abstract Syntax Tree}) es la representaci\'on intermedia est\'andar
en los compiladores e int\'erpretes.  Representa la estructura sint\'actica
de un programa mediante un \'arbol cuyos nodos corresponden a los constructores
del lenguaje y cuyos hijos representan los subcomponentes de cada
constructor~\cite{aho2006}.

El AST es \emph{abstracto} en el sentido de que omite los detalles de la
sintaxis concreta ---par\'entesis, palabras clave, signos de puntuaci\'on---
y conserva \'unicamente la estructura l\'ogica del programa~\cite{scott2015}.

En el dise\~no del DSL presentado en este trabajo, el AST cumple el papel
de \textbf{representaci\'on intermedia independiente del backend}: el
programador construye el AST mediante las formas del DSL, y el backend
lo recorre para producir el c\'odigo Python de generaci\'on del \xlsx{}.
Gracias a esta separaci\'on, el mismo AST podr\'{\i}a en principio alimentar
backends alternativos (LibreOffice Calc, CSV, HTML) sin modificar la capa
del DSL.

Cada nodo del AST es una instancia de una clase CLOS definida con la macro
\dsl{defclass*}.  Los nodos se dividen en cuatro familias:
\begin{enumerate}
  \item Nodos de estructura del workbook (\dsl{xl-workbook}, \dsl{xl-sheet},
    \dsl{xl-region}, \dsl{xl-table}).
  \item Nodos de expresi\'on de celda (\dsl{xl-expr-column-ref},
    \dsl{xl-table-range}, \dsl{xl-expr-countif}, \dsl{xl-expr-if}, etc.).
  \item Nodos de cuantificaci\'on existencial (\dsl{xl-expr-exists},
    \dsl{xl-domain-rows}).
  \item Nodos de formato condicional y posicionamiento (\dsl{xl-style-rule},
    \dsl{xl-anchor}, \dsl{xl-nav}, \dsl{xl-sheet-fixed-expr}).
\end{enumerate}

Esta organizaci\'on se describe en detalle en el
cap\'{\i}tulo~\ref{cap:dsl}.

% ---------------------------------------------------------------------------
\section{El lenguaje de programaci\'on Common Lisp}
\label{sec:lisp}
% ---------------------------------------------------------------------------

Common Lisp es un lenguaje de programaci\'on de prop\'osito general creado
originalmente por John McCarthy en 1958~\cite{mccarthy1960} y estandarizado
como Common Lisp en 1994~\cite{steele1990}.  Su elecci\'on como plataforma
de implementaci\'on del DSL se justifica por cuatro caracter\'{\i}sticas
decisivas.

% ............................................................................
\subsection{Notaci\'on prefija y homoiconicidad}
\label{ssec:prefija}
% ............................................................................

En Common Lisp, todas las expresiones son listas anidadas con notaci\'on
prefija: el primer elemento es el operador y los siguientes son los
argumentos.  Una llamada a funci\'on, la definici\'on de una clase, una
iteraci\'on y una expresi\'on del DSL tienen exactamente la misma forma
sint\'actica.

Esta propiedad, llamada \textbf{homoiconicidad}, hace que el c\'odigo Lisp
sea indistinguible de los datos Lisp~\cite{graham1993}.  El compilador puede
inspeccionar y transformar el c\'odigo del mismo modo que inspecciona y
transforma listas de datos.  Los DSL internos se benefician directamente
de esto: las formas del DSL son listas ordinarias que los macros transforman
en c\'odigo v\'alido.

% ............................................................................
\subsection{Macros}
\label{ssec:macros}
% ............................................................................

Un \textbf{macro} de Lisp es una funci\'on que opera sobre el c\'odigo fuente
del programa ---representado como listas--- y devuelve c\'odigo nuevo que
sustituye a la forma original antes de la compilaci\'on o evaluaci\'on~\cite{seibel2005}.

La diferencia con una funci\'on ordinaria es temporal: una funci\'on se
eval\'ua en tiempo de ejecuci\'on; un macro se expande en tiempo de
compilaci\'on.  Esta distinci\'on permite a los macros introducir nueva
sintaxis y nuevas abstracciones sin coste en tiempo de ejecuci\'on.

En este trabajo, cada constructo del DSL ---\dsl{def-table},
\dsl{conditional-rendering}, \dsl{exists}, \dsl{nav}--- se implementa como
un macro que construye uno o m\'as nodos del AST.  El programador escribe en
vocabulario del dominio; el macro transforma esa escritura en llamadas a los
constructores CLOS.

% ............................................................................
\subsection{El Sistema de Objetos de Common Lisp (CLOS)}
\label{ssec:clos}
% ............................................................................

El \textbf{Sistema de Objetos de Common Lisp} (CLOS, del ingl\'es
\textit{Common Lisp Object System}) es el sistema de orientaci\'on a objetos
integrado en el est\'andar del lenguaje~\cite{clos1988}.  Sus caracter\'{\i}sticas
principales son:

\begin{description}
  \item[Herencia m\'ultiple.] Una clase puede heredar de varias superclases
    simult\'aneamente.  Los conflictos de herencia se resuelven mediante el
    \textit{Class Precedence List} (CPL), un orden lineal determinista~\cite{keene1989}.

  \item[Despacho m\'ultiple.] Los m\'etodos no pertenecen a una clase
    espec\'{\i}fica.  Una funci\'on gen\'erica puede despachar sobre los tipos
    de varios argumentos simult\'aneamente~\cite{kiczales1991}.  Esto significa
    que el c\'omportamiento de \dsl{(compile-excel-formula nodo ...)} depende
    del tipo de \dsl{nodo}, no de un m\'etodo ligado a una clase particular.

  \item[Funciones gen\'ericas.] Las funciones polim\'orficas se declaran con
    \dsl{defgeneric} y se especializan con \dsl{defmethod}.  A\~nadir soporte
    para un nuevo tipo de nodo AST s\'olo requiere a\~nadir un nuevo
    \dsl{defmethod}; el resto del sistema no se modifica.
\end{description}

En este trabajo, CLOS cumple dos funciones.  Primero, las clases CLOS
representan los nodos del AST: cada tipo de nodo es una clase con slots
que almacenan sus atributos.  Segundo, la funci\'on gen\'erica
\dsl{compile-excel-formula} usa el despacho m\'ultiple para compilar
cada tipo de nodo a su cadena de f\'ormula Excel correspondiente, sin
ninguna estructura \dsl{cond} o \dsl{case} centralizada.

% ............................................................................
\subsection{La macro \texttt{defclass*}}
\label{ssec:defclass-star}
% ............................................................................

Para reducir el c\'odigo repetitivo de definir clases CLOS y sus
constructores, el sistema define la macro \dsl{defclass*}, que genera
simult\'aneamente:

\begin{itemize}
  \item Una clase CLOS con nombre \dsl{clase-xl-\textit{nombre}}, donde
    cada slot tiene un accessor y un \texttt{initarg} del mismo nombre.
  \item Una funci\'on constructora \dsl{xl-\textit{nombre}} con
    argumentos por clave (\dsl{\&key}), que crea una instancia de la clase
    con los valores proporcionados.
\end{itemize}

\begin{lstlisting}[language=CommonLisp,
  caption={Implementaci\'on de defclass*.}]
(defmacro defclass* (class-name parents slots
                      &key (ctr-args-type '&key))
  `(progn
     (defclass ,(symb 'clase- class-name) ,parents
       ,(loop for s in slots
              collecting (expansion-del-slot s)))
     ,(make-constructor class-name slots
                        :argument-type ctr-args-type)))
\end{lstlisting}

Por ejemplo, la definici\'on del nodo \dsl{xl-expr-if}:

\begin{lstlisting}[language=CommonLisp,
  caption={Ejemplo de uso de defclass*.}]
(defclass* xl-expr-if () (test then else))
;; genera:
;;   (defclass clase-xl-expr-if () (test :accessor test :initarg :test)
;;                                  (then :accessor then :initarg :then)
;;                                  (else :accessor else :initarg :else))
;;   (defun xl-expr-if (&key test then else)
;;     (make-instance 'clase-xl-expr-if :test test :then then :else else))
\end{lstlisting}

Este patr\'on elimina el c\'odigo repetitivo de definici\'on de nodos y
garantiza que cada clase tiene una funci\'on constructora consistente con
su interfaz de inicializaci\'on.

% ===========================================================================
%  Capítulo 2 — Análisis del dominio de planificación tabular
% ===========================================================================
\chapter{An\'alisis del dominio de planificaci\'on tabular}
\label{cap:dominio}

Para dise\~nar un lenguaje que capture los patrones recurrentes de la
planificaci\'on tabular, es necesario estudiar primero los casos concretos
que el lenguaje debe resolver.  En este cap\'{\i}tulo se presentan tres
casos de uso reales de la Facultad de Matem\'atica y Computaci\'on de la
Universidad de La Habana, se identifican los elementos comunes entre ellos
y se derivan los requisitos del lenguaje.

% ---------------------------------------------------------------------------
\section{Casos de uso}
\label{sec:casos}
% ---------------------------------------------------------------------------

% ............................................................................
\subsection{Horario de defensas de tesis}
\label{ssec:caso-defensas}
% ............................................................................

El proceso de planificaci\'on de defensas de tesis de diploma asigna a cada
estudiante un conjunto de cinco profesores que forman el tribunal ---tutor,
oponente, presidente, vocal y secretario---, una fecha y hora de exposici\'on,
y un local donde esta se celebra.  El coordinador de defensas debe detectar
conflictos: un mismo profesor asignado a dos defensas con el mismo slot
horario (mismo d\'{\i}a y hora), lo que hace imposible su presencia en
ambas.

La hoja de c\'alculo que da soporte a este proceso contiene dos tablas:

\begin{description}
  \item[Tabla de defensas.] Una fila por defensa, con columnas para el
    estudiante y los cinco roles del tribunal, m\'as las columnas de d\'{\i}a,
    hora y local.  Esta tabla no necesita columnas calculadas: toda la
    informaci\'on se introduce manualmente.

  \item[Tabla de profesores.] Una fila por profesor, con su nombre, grado
    acad\'emico y el n\'umero de defensas en que participa.  Este \'ultimo dato
    se calcula autom\'aticamente contando cu\'antas veces aparece el nombre del
    profesor en cualquiera de los cinco roles de la tabla de defensas.
    Adem\'as, las filas de los profesores con conflicto de horario se resaltan
    en rojo.
\end{description}

Los datos del caso de validaci\'on incluyen tres defensas y seis profesores.
Dos de las defensas comparten el slot \texttt{Lunes~10:00} con integrantes
de tribunal en com\'un, lo que genera un conflicto que la hoja debe detectar
y resaltar:

\begin{table}[h!]
\centering
\caption{Datos de las defensas del caso de validaci\'on.}
\label{tab:datos-defensas}
\renewcommand{\arraystretch}{1.2}
\begin{tabular}{llllllll}
\hline
\textbf{Estudiante} & \textbf{Tutor} & \textbf{Opon.}
  & \textbf{Pres.} & \textbf{Vocal} & \textbf{Secr.}
  & \textbf{D\'{\i}a} & \textbf{Hora} \\
\hline
Juan D\'{\i}az  & Piad   & Garc\'{\i}a & L\'opez  & Torres & Ruiz   & Lunes & 10:00 \\
Mar\'{\i}a Vega & Garc\'{\i}a & L\'opez  & Piad   & Soto   & Torres & Lunes & 14:00 \\
Luis Mora   & L\'opez  & Torres & Garc\'{\i}a & Ruiz   & Soto   & Lunes & 10:00 \\
\hline
\end{tabular}
\end{table}

Las defensas de Juan D\'{\i}az y Luis Mora comparten el slot
\texttt{Lunes~10:00}.  El profesor L\'opez participa como oponente en la
primera y como tutor en la segunda; el profesor Torres como vocal en la
primera y como oponente en la segunda.  Ambos tienen conflicto y deben
aparecer resaltados en rojo en la tabla de profesores.

Este caso es el que se usa como ejemplo motivador a lo largo del cap\'{\i}tulo~\ref{cap:dsl}.

% ............................................................................
\subsection{Horario docente de la facultad}
\label{ssec:caso-facultad}
% ............................................................................

La asignaci\'on de docencia en la Facultad de Matem\'atica y Computaci\'on
implica distribuir los turnos de clase de cada semana entre los grupos de las
tres carreras (Ciencia de la Computaci\'on, Matem\'atica y Ciencia de Datos),
en cuatro a\~nos cada una.  Cada grupo tiene una hoja propia en el workbook
con su horario semanal.

La hoja de c\'alculo de planificaci\'on contiene, para cada grupo, tres tablas
dispuestas en la misma fila de la hoja:

\begin{description}
  \item[Tabla de turnos (\texttt{turno-table}).] Seis columnas (turno,
    lunes, martes, mi\'ercoles, jueves y viernes), con dos filas Excel por
    turno l\'ogico (\dsl{cell-height~=~2}): la primera subfila almacena el
    nombre de la asignatura y la segunda el aula asignada.

  \item[Tabla de estad\'{\i}sticas (\texttt{stats-table}).] Una fila por
    asignatura, con su abreviatura, nombre completo, frecuencia semanal
    requerida, asignaciones ya realizadas (calculadas contando las apariciones
    de la abreviatura en la tabla de turnos) y horas que faltan por asignar
    (calculadas como la diferencia entre frecuencia y asignadas).

  \item[Tabla de aulas (\texttt{aulas-table}).] Lista de aulas disponibles
    para el grupo.
\end{description}

Adem\'as de las hojas de grupo, existe una hoja de ocupaci\'on de aulas que
agrega, para cada turno de cada d\'{\i}a, los grupos que ocupan cada aula.
Esta tabla usa referencias cruzadas din\'amicas (\dsl{collect-over}) que
iteran sobre todas las hojas de grupo y extraen el aula asignada en cada
turno.  Diecisiete grupos generan diecisiete hojas y una hoja de aulas con
referencias a las diecisiete.

Este caso pone a prueba dos caracter\'{\i}sticas del DSL que el caso de
defensas no ejercita: las tablas con celdas de altura compuesta
(\dsl{cell-height} mayor que 1) y las referencias cruzadas entre hojas del
mismo workbook mediante \dsl{collect-over}.

% ............................................................................
\subsection{Horario de programaci\'on televisiva}
\label{ssec:caso-tv}
% ............................................................................

El tercer caso de uso es el horario de programaci\'on semanal de un canal de
televisi\'on.  Cada programa tiene una hora de inicio, una duraci\'on en
minutos, un nombre, un tipo y un p\'ublico objetivo.  La hoja de c\'alculo
muestra una p\'agina por d\'{\i}a, con columnas para cada uno de esos campos
m\'as una columna de hora de fin, calculada autom\'aticamente a partir de la
hora de inicio y la duraci\'on.

La f\'ormula de la hora de fin involucra operaciones de tiempo:
\[
\texttt{fin} = \texttt{TEXT}(\texttt{TIMEVALUE}(\texttt{inicio}) + \texttt{duraci\'on}/1440, \texttt{"hh:mm"})
\]
donde $1440$ es el n\'umero de minutos en un d\'{\i}a.  Adem\'as, cada fila
se colorea seg\'un el tipo de programa (informativo, musical, infantil, etc.)
mediante un color distinto en la celda.

Este caso valida que el DSL puede expresar operaciones aritm\'eticas sobre
valores de tiempo (\dsl{time-add}) y f\'ormulas de celda simples sin
cuantificadores existenciales.

% ---------------------------------------------------------------------------
\section{Elementos comunes}
\label{sec:elementos-comunes}
% ---------------------------------------------------------------------------

Del an\'alisis de los tres casos se identifican cuatro patrones que se
repiten en todos o en la mayor\'{\i}a de los documentos de planificaci\'on
tabular:

% ............................................................................
\subsection{Tablas relacionadas y referencias cruzadas}
\label{ssec:elem-tablas}
% ............................................................................

En todos los casos, el documento contiene m\'as de una tabla cuyos datos est\'an
relacionados.  La tabla de profesores refencia la tabla de defensas; la tabla
de estad\'{\i}sticas referencia la tabla de turnos; la hoja de aulas referencia
las hojas de cada grupo.

Estas referencias son de dos tipos:
\begin{itemize}
  \item \textbf{Intrarregi\'on:} ambas tablas est\'an en la misma hoja y
    se accede a ellas mediante rangos absolutos de columnas (p.~ej.,
    \texttt{\$B\$4:\$F\$6} para los cinco roles de la tabla de defensas).
  \item \textbf{Interhoja:} la tabla referenciada est\'a en otra hoja del
    mismo workbook (p.~ej., la hoja de aulas referencia las hojas de grupos).
\end{itemize}

El lenguaje debe resolver estas referencias sin que el programador escriba
letras de columna ni nombres de hoja: los identifica por los s\'{\i}mbolos
DSL de las columnas y por el identificador de la tabla.

% ............................................................................
\subsection{Columnas calculadas}
\label{ssec:elem-computed}
% ............................................................................

En todos los casos, al menos una columna de alguna tabla muestra un valor
calculado a partir de datos de otra tabla.  Los tipos de c\'alculo m\'as
frecuentes son:

\begin{description}
  \item[\texttt{COUNTIF}.] Cuenta cu\'antas celdas de un rango cumplen un
    criterio.  Ejemplo: n\'umero de veces que aparece el nombre de un
    profesor en los cinco roles de la tabla de defensas.

  \item[\texttt{COUNTA}.] Cuenta celdas no vac\'{\i}as en un rango.  Ejemplo:
    n\'umero de defensas registradas.

  \item[\texttt{SUM}.] Suma un rango.  Ejemplo: total de defensas acumuladas
    por todos los profesores.

  \item[Aritmética de tiempo.] Suma una duraci\'on en minutos a una hora de
    inicio.  Ejemplo: hora de fin de un programa televisivo.

  \item[\texttt{collect-over}.] Agrega resultados de la misma posici\'on en
    m\'ultiples hojas.  Ejemplo: grupos que ocupan el Aula~1 en el turno~3
    del lunes.
\end{description}

% ............................................................................
\subsection{Visualizaci\'on condicional}
\label{ssec:elem-condicional}
% ............................................................................

En los casos de defensas y de docencia, la condici\'on de coloreado depende
de datos de otra tabla.  El color no es est\'atico: si el usuario edita los
datos directamente en Excel, el coloreado debe recalcularse autom\'aticamente.

El mecanismo adecuado es, por tanto, una \texttt{FormulaRule} de formato
condicional incrustada en el archivo \xlsx{}, no un coloreado est\'atico
aplicado en tiempo de generaci\'on.  La condici\'on de la regla es una
f\'ormula SUMPRODUCT que captura la sem\'antica del cuantificador existencial
del DSL.

La forma en que el DSL expresa esta condici\'on ---mediante \dsl{exists}---
y la forma en que el backend la implementa ---mediante SUMPRODUCT--- se
detallan en los cap\'{\i}tulos~\ref{cap:dsl} y~\ref{cap:generacion},
respectivamente.

% ............................................................................
\subsection{Posicionamiento relativo}
\label{ssec:elem-posicion}
% ............................................................................

Todos los casos requieren expresiones fuera de las tablas: totales, etiquetas,
celdas de resumen.  La posici\'on de estas expresiones depende de cu\'antas
filas de datos tiene la tabla, que var\'{\i}a en cada ejecuci\'on.

La soluci\'on adoptada en el DSL es el sistema de anclas y navegaci\'on
(\dsl{nav}, \dsl{ultima-fila}, \dsl{primera-fila}): se declara que una
expresi\'on va ``una fila por debajo de la \'ultima fila de datos de la
columna \texttt{estudiante} de la tabla de defensas'', y el compilador
resuelve esa declaraci\'on a la coordenada Excel concreta en tiempo de
generaci\'on.

% ---------------------------------------------------------------------------
\section{Requisitos del lenguaje}
\label{sec:requisitos}
% ---------------------------------------------------------------------------

Del an\'alisis anterior se derivan los requisitos que el DSL debe satisfacer:

\begin{enumerate}
  \item \textbf{Declaraci\'on de tablas.} El lenguaje debe permitir definir
    tablas con nombre, con una lista de columnas identificadas por s\'{\i}mbolos,
    y con soporte para celdas de altura y anchura compuesta.

  \item \textbf{Columnas computadas.} Debe ser posible asociar una expresi\'on
    del DSL a una columna; el compilador genera la f\'ormula Excel
    correspondiente para cada fila de datos.

  \item \textbf{Referencias intrarregi\'on.} El lenguaje debe resolver
    referencias a columnas de la misma tabla (\dsl{col}) y a rangos de
    columnas de otra tabla en la misma regi\'on (\dsl{trange}, \dsl{tcol}).

  \item \textbf{Cuantificaci\'on existencial.} El lenguaje debe ofrecer una
    forma de expresar condiciones de existencia sobre filas de la misma tabla
    (\dsl{:all-rows}) o de otra tabla (\dsl{:rows-of}), especificando
    columnas de coincidencia y columnas de solapamiento.

  \item \textbf{Formato condicional persistente.} Las reglas de coloreado
    condicional deben traducirse a \texttt{FormulaRule} en el \xlsx{}, de
    modo que persistan y se reevaluen cuando el usuario edita el archivo.

  \item \textbf{Posicionamiento relativo.} Debe ser posible colocar
    expresiones fijas en posiciones calculadas como un ancla m\'as pasos
    direccionales, sin coordenadas absolutas codificadas.

  \item \textbf{Referencias interhoja.} El lenguaje debe soportar referencias
    din\'amicas a celdas de otras hojas mediante iteraci\'on sobre una lista
    de nombres de hojas (\dsl{collect-over}).

  \item \textbf{Ensamblaje declarativo.} El workbook completo debe poderse
    describir con formas del DSL (\dsl{libro}, \dsl{hoja}, \dsl{tabla})
    sin ning\'un conocimiento de las coordenadas del archivo de salida.

  \item \textbf{Independencia del backend.} El AST no debe contener
    ning\'un detalle del target de generaci\'on.  Agregar un nuevo tipo de
    expresi\'on no debe requerir modificar el DSL; solo a\~nadir un nuevo
    m\'etodo al backend.
\end{enumerate}

La tabla~\ref{tab:requisitos-casos} muestra c\'omo cada requisito se
manifiesta en los casos de uso analizados.

\begin{table}[h!]
\centering
\caption{Requisitos del lenguaje y casos de uso en que se manifiestan.}
\label{tab:requisitos-casos}
\renewcommand{\arraystretch}{1.2}
\begin{tabular}{p{5.2cm} ccc}
\toprule
\textbf{Requisito} & \textbf{Defensas} & \textbf{Facultad} & \textbf{TV} \\
\midrule
Declaraci\'on de tablas         & \checkmark & \checkmark & \checkmark \\
Columnas computadas              & \checkmark & \checkmark & \checkmark \\
Referencias intrarregi\'on       & \checkmark & \checkmark & --- \\
Cuantificaci\'on existencial     & \checkmark & --- & --- \\
Formato condicional persistente  & \checkmark & --- & --- \\
Posicionamiento relativo         & \checkmark & \checkmark & --- \\
Referencias interhoja            & --- & \checkmark & --- \\
Ensamblaje declarativo           & \checkmark & \checkmark & \checkmark \\
Independencia del backend        & \checkmark & \checkmark & \checkmark \\
\bottomrule
\end{tabular}
\end{table}

% ===========================================================================
%  Capítulo — El lenguaje de dominio específico
% ===========================================================================
\chapter{El Lenguaje de Dominio Espec\'{\i}fico}
\label{cap:dsl}

% ---------------------------------------------------------------------------
\section{Visi\'on general}
\label{sec:vision-general}
% ---------------------------------------------------------------------------

Numerosos problemas cotidianos producen informaci\'on que se organiza
naturalmente en forma de tabla: filas que representan entidades y columnas
que describen sus atributos, relaciones o restricciones.  Los horarios son
un ejemplo paradigm\'atico de esta clase de datos---franjas horarias,
participantes, recursos y conflictos pueden expresarse como tablas
relacionadas---pero la misma estructura aparece en inventarios, resultados
acad\'emicos, planificaci\'on de proyectos y cualquier dominio donde la
informaci\'on admita una organizaci\'on tabular.

El objetivo del lenguaje de dominio espec\'{\i}fico (DSL) desarrollado en
este trabajo es ofrecer un recurso de descripci\'on para esa clase de datos:
un vocabulario declarativo con el que el programador expresa \emph{qu\'e}
contiene cada tabla---columnas, f\'ormulas derivadas, relaciones entre
tablas, reglas de presentaci\'on---sin ocuparse de \emph{c\'omo} se
materializa esa descripci\'on en un formato concreto.

El DSL est\'a implementado como un conjunto de funciones y macros de
Common~Lisp que construyen, en tiempo de carga, un \emph{\'arbol de sintaxis
abstracta} (AST) cuyos nodos son instancias de clases CLOS.  La
canalizaci\'on de compilaci\'on comprende tres etapas claramente separadas:

\begin{enumerate}
  \item \textbf{Expansi\'on de macros.}  Cada forma del DSL
    (p.\,ej.\ \texttt{def-table}, \texttt{exists}) se expande en una llamada
    al constructor CLOS del nodo AST correspondiente.  El resultado es un
    \'arbol cuyos nodos representan las entidades a partir de las cuales el
    sistema es capaz de generar las representaciones visuales de los datos.

  \item \textbf{Generaci\'on de c\'odigo.}  El \'arbol se recorre por el
    backend, que traduce cada nodo a las instrucciones propias del formato de
    salida: expresiones de celda, reglas de presentaci\'on y directivas de
    construcci\'on del documento.  El backend es el \'unico punto que conoce
    los nombres de columna reales, los \'{\i}ndices de fila y los detalles
    de la API de destino.

  \item \textbf{Ejecuci\'on.}  El programa generado se ejecuta, produciendo
    el archivo de salida con los datos, las f\'ormulas derivadas y las reglas
    de presentaci\'on incrustadas.
\end{enumerate}

Un principio central de dise\~no es la \emph{separaci\'on entre intenci\'on
e implementaci\'on}: el DSL expresa la sem\'antica del problema
(``resaltar los profesores que tienen un conflicto de horario'') y el backend
decide el mecanismo concreto seg\'un el formato de destino elegido.  El
programador del DSL no escribe ninguna instrucci\'on espec\'{\i}fica del
formato de salida.

% ---------------------------------------------------------------------------
\section{\'Arbol de sintaxis abstracta}
\label{sec:ast-nodos}
% ---------------------------------------------------------------------------

El \'arbol de sintaxis abstracta (AST) es la representaci\'on intermedia que
separa la capa del DSL de la capa del generador de c\'odigo.  Cada concepto
del problema---una tabla, una expresi\'on de celda, un cuantificador
existencial---queda capturado como un nodo AST con atributos propios.  El
backend opera \'unicamente sobre este \'arbol; no interpreta texto ni
inspecciona las formas del DSL.

Todos los nodos del AST se definen con la macro \texttt{defclass*}, que
genera simult\'aneamente una clase CLOS \texttt{clase-xl-\textit{nombre}} y
una funci\'on constructora \texttt{xl-\textit{nombre}} con par\'ametros por
clave.  Por ejemplo, el nodo del cuantificador existencial se define y
construye de la siguiente manera:

\begin{lstlisting}[language=CommonLisp,
  caption={Patr\'on de definici\'on de nodo AST con defclass*.}]
;; defclass* genera la clase y el constructor xl-expr-exists:
(defclass* xl-expr-exists ()
           (bind-var domain match-keys overlap-cols self-in-cols))

;; El constructor acepta los slots como :keywords:
(xl-expr-exists :bind-var     'r
                :domain       (xl-domain-rows :table-id nil)
                :match-keys   '(dia hora)
                :overlap-cols '(tutor oponente presidente vocal secretario)
                :self-in-cols  nil)
\end{lstlisting}

Los nodos del AST se agrupan en seis categor\'{\i}as seg\'un su papel en
el \'arbol: estructura del workbook, expresiones de celda, navegaci\'on de
fila, cuantificaci\'on existencial, expresiones cross-sheet din\'amicas y
formato condicional con posicionamiento relativo.

% ............................................................................
\subsection{Estructura del workbook}
\label{sec:ast-estructura}
% ............................................................................

Los nodos de estructura representan los contenedores del libro.
Son los nodos de m\'as alto nivel del \'arbol: el workbook agrupa hojas, las
hojas agrupan regiones y las regiones agrupan tablas.  Su jerarqu\'{\i}a
refleja directamente la estructura del documento de salida.

\begin{table}[h!]
\centering
\caption{Nodos de estructura del workbook.}
\label{tab:ast-estructura}
\renewcommand{\arraystretch}{1.3}
\begin{tabular}{p{3.6cm} p{4.8cm} p{5.4cm}}
\toprule
\textbf{Nodo} & \textbf{Atributos} & \textbf{Papel} \\
\midrule
\texttt{xl-workbook}
  & \texttt{name}, \texttt{sheets}
  & Ra\'{\i}z del \'arbol.  \texttt{name} es el nombre del
    archivo de salida; \texttt{sheets} es la lista de hojas. \\
\midrule
\texttt{xl-sheet}
  & \texttt{name}, \texttt{regions},\newline
    \texttt{fixed-expressions}
  & Una hoja del libro.  Contiene regiones horizontales y
    expresiones en posici\'on relativa fuera de las tablas. \\
\midrule
\texttt{xl-region}
  & \texttt{tables}
  & Agrupa tablas colocadas una junto a otra en la misma
    hoja, separadas por una columna en blanco. \\
\midrule
\texttt{xl-table}
  & \texttt{id}, \texttt{col-names},\newline
    \texttt{headers}, \texttt{contenido},\newline
    \texttt{computed}, \texttt{style-rules},\newline
    \texttt{cell-height}, \texttt{cell-width}
  & Plantilla de tabla.  \texttt{col-names} enumera los
    s\'{\i}mbolos DSL de columna; \texttt{computed} asocia
    columnas a expresiones de f\'ormula; \texttt{style-rules}
    almacena las reglas de coloreado condicional. \\
\bottomrule
\end{tabular}
\end{table}

% ............................................................................
\subsection{Expresiones de celda}
\label{sec:ast-expresiones}
% ............................................................................

Las expresiones de celda son los nodos que el backend traduce a
representaciones de f\'ormula en el formato de salida.  Abarcan tanto
referencias a columnas de la misma tabla o de otra (\texttt{xl-expr-column-ref},
\texttt{xl-table-range}) como operaciones de agregaci\'on (\texttt{COUNTIF},
\texttt{COUNTA}, \texttt{SUM}) y literales de texto.  Gracias a esta
representaci\'on, el programador escribe \texttt{(countif rango criterio)}
en lugar de componer manualmente la cadena de f\'ormula con coordenadas
absolutas.

\begin{table}[h!]
\centering
\caption{Nodos de expresi\'on de celda.}
\label{tab:ast-expresiones}
\renewcommand{\arraystretch}{1.3}
\begin{tabular}{p{4.2cm} p{3.8cm} p{5.8cm}}
\toprule
\textbf{Nodo} & \textbf{Atributos} & \textbf{Papel} \\
\midrule
\texttt{xl-expr-column-ref}
  & \texttt{name}, \texttt{context}
  & Referencia a la columna \texttt{name} en la fila actual.
    Se resuelve a la coordenada de esa columna en el formato de salida. \\
\midrule
\texttt{xl-range}
  & \texttt{from-col}, \texttt{to-col}
  & Rango de columnas contiguas de la tabla actual. \\
\midrule
\texttt{xl-table-range}
  & \texttt{table-id},\newline
    \texttt{from-col}, \texttt{to-col}
  & Rango en una tabla espec\'{\i}fica de la regi\'on
    (referencia cross-table).  Lo produce la forma
    \texttt{trange}. \\
\midrule
\texttt{xl-expr-countif}
  & \texttt{count-range},\newline\texttt{criteria}
  & F\'ormula de conteo condicional sobre un rango. \\
\midrule
\texttt{xl-expr-counta}
  & \texttt{count-range}
  & Conteo de celdas no vac\'{\i}as. \\
\midrule
\texttt{xl-expr-sum}
  & \texttt{count-range}
  & Suma sobre un rango. \\
\midrule
\texttt{xl-expr-string}
  & \texttt{value}
  & Literal de texto; produce una cadena fija en la celda. \\
\midrule
\texttt{xl-expr-if}
  & \texttt{test}, \texttt{then}, \texttt{else}
  & Condicional de tres ramas. \\
\midrule
\texttt{xl-expr-equals}
  & \texttt{a}, \texttt{b}
  & Igualdad entre dos expresiones. \\
\midrule
\texttt{xl-expr-and}, \texttt{xl-expr-or}
  & \texttt{a}, \texttt{b}
  & Conjunci\'on y disyunci\'on l\'ogica. \\
\midrule
\texttt{xl-expr-concat}
  & \texttt{a}, \texttt{b}
  & Concatenaci\'on de texto. \\
\midrule
\texttt{xl-expr-table-col-ref}
  & \texttt{table-id}, \texttt{name},\newline\texttt{context}
  & Referencia a una columna de una tabla espec\'{\i}fica de la regi\'on.
    Evita la ambig\"uedad cuando dos tablas comparten el mismo nombre de
    columna. \\
\midrule
\texttt{xl-expr-param-ref}
  & \texttt{name}
  & Referencia a un par\'ametro de instancia almacenado en celda auxiliar. \\
\midrule
\texttt{xl-expr-different}
  & \texttt{a}, \texttt{b}
  & Desigualdad entre dos expresiones. \\
\midrule
\texttt{xl-expr-gt}, \texttt{xl-expr-lt}
  & \texttt{a}, \texttt{b}
  & Comparaci\'on estricta: mayor que y menor que. \\
\midrule
\texttt{xl-expr-gte}, \texttt{xl-expr-lte}
  & \texttt{a}, \texttt{b}
  & Comparaci\'on no estricta: mayor o igual y menor o igual. \\
\midrule
\texttt{xl-expr-add}, \texttt{xl-expr-subtract}
  & \texttt{a}, \texttt{b}
  & Suma y resta aritm\'etica. \\
\midrule
\texttt{xl-expr-multiply}, \texttt{xl-expr-divide}
  & \texttt{a}, \texttt{b}
  & Multiplicaci\'on y divisi\'on aritm\'etica. \\
\midrule
\texttt{xl-expr-time-add}
  & \texttt{a}, \texttt{b}
  & Suma de tiempos como texto:
    \texttt{TEXT(TIMEVALUE(a)+(b/1440),"hh:mm")}. \\
\midrule
\texttt{xl-expr-lookup}
  & \texttt{value-field},\newline\texttt{key-expr}
  & B\'usqueda en la tabla actual: devuelve la columna
    \texttt{value-field} correspondiente a \texttt{key-expr}. \\
\bottomrule
\end{tabular}
\end{table}

% ............................................................................
\subsection{Navegaci\'on de fila y control de visibilidad}
\label{sec:ast-navegacion-fila}
% ............................................................................

Ciertos c\'alculos necesitan referenciar la fila anterior o siguiente de la
misma tabla---por ejemplo, acumular un valor incremental o mostrar una
transici\'on---o bien producir una celda visualmente vac\'{\i}a bajo alguna
condici\'on.  Estos nodos encapsulan esa l\'ogica sin exponer \'{\i}ndices
de fila al programador del DSL.

\begin{table}[h!]
\centering
\caption{Nodos de navegaci\'on de fila y control de visibilidad.}
\label{tab:ast-navegacion}
\renewcommand{\arraystretch}{1.3}
\begin{tabular}{p{3.8cm} p{3.8cm} p{6.2cm}}
\toprule
\textbf{Nodo} & \textbf{Atributos} & \textbf{Papel} \\
\midrule
\texttt{xl-expr-previous-row}
  & \texttt{expr}
  & Eval\'ua \texttt{expr} referenciando la columna en la
    fila anterior de la tabla. \\
\midrule
\texttt{xl-expr-next-row}
  & \texttt{expr}
  & Eval\'ua \texttt{expr} referenciando la columna en la
    fila siguiente de la tabla. \\
\midrule
\texttt{xl-expr-first-row}
  & ---
  & Produce \texttt{TRUE} cuando la fila en evaluaci\'on es la
    primera fila de datos. \\
\midrule
\texttt{xl-expr-non-empty}
  & \texttt{expr}
  & Produce \texttt{TRUE} cuando \texttt{expr} no es vac\'{\i}a
    ni cero. \\
\midrule
\texttt{xl-expr-show-nothing}
  & ---
  & Produce una cadena vac\'{\i}a: la celda queda visualmente
    vac\'{\i}a. \\
\bottomrule
\end{tabular}
\end{table}

% ............................................................................
\subsection{Cuantificaci\'on existencial}
\label{sec:ast-cuantificacion}
% ............................................................................

La cuantificaci\'on existencial es el mecanismo del DSL para expresar
condiciones de la forma ``existe alguna fila que cumpla ciertos criterios''.
Mantenerla como un nodo AST propio, y no como una expresi\'on booleana
expl\'{\i}cita, permite al backend generar la comprobaci\'on m\'as adecuada
seg\'un el contexto: una f\'ormula SUMPRODUCT vectorial en el backend Excel,
u otra implementaci\'on para un backend diferente.

\begin{table}[h!]
\centering
\caption{Nodos de cuantificaci\'on existencial.}
\label{tab:ast-cuantificacion}
\renewcommand{\arraystretch}{1.3}
\begin{tabular}{p{3.6cm} p{4.8cm} p{5.4cm}}
\toprule
\textbf{Nodo} & \textbf{Atributos} & \textbf{Papel} \\
\midrule
\texttt{xl-expr-exists}
  & \texttt{bind-var}, \texttt{domain},\newline
    \texttt{match-keys},\newline
    \texttt{overlap-cols},\newline
    \texttt{self-in-cols}
  & Cuantificador existencial estructurado.  Declara que debe
    existir una fila del dominio que cumpla las condiciones
    expresadas por los slots sem\'anticos.  El backend decide
    el mecanismo de comprobaci\'on seg\'un el tipo de dominio. \\
\midrule
\texttt{xl-domain-rows}
  & \texttt{table-id}
  & Dominio de filas.  \texttt{table-id = nil} indica la
    tabla actual; un s\'{\i}mbolo identifica una tabla
    espec\'{\i}fica de la regi\'on (dominio cross-table). \\
\bottomrule
\end{tabular}
\end{table}

% ............................................................................
\subsection{Expresiones cross-sheet din\'amicas}
\label{sec:ast-cross-sheet}
% ............................................................................

Cuando un libro contiene m\'ultiples hojas con estructura homog\'enea---por
ejemplo, una hoja por grupo o aula---es necesario agregar valores a trav\'es
de ellas.  A diferencia de \texttt{xl-expr-cross-sheet-ref}, que referencia
una hoja con nombre fijo, los nodos de esta secci\'on trabajan con una
\emph{variable de hoja} resuelta en tiempo de compilaci\'on mediante el
entorno din\'amico \texttt{*sheet-env*}, declarado por \texttt{collect-over}.

\begin{table}[h!]
\centering
\caption{Nodos de expresi\'on cross-sheet din\'amica.}
\label{tab:ast-cross-sheet}
\renewcommand{\arraystretch}{1.3}
\begin{tabular}{p{3.8cm} p{4.4cm} p{5.6cm}}
\toprule
\textbf{Nodo} & \textbf{Atributos} & \textbf{Papel} \\
\midrule
\texttt{xl-expr-cross-sheet-ref}
  & \texttt{sheet},\newline\texttt{cell-template}
  & Referencia a una celda en otra hoja con nombre fijo.
    \texttt{cell-template} puede incluir \texttt{\$\{row-num\}} para
    fila din\'amica. \\
\midrule
\texttt{xl-expr-cross-cell}
  & \texttt{sheet}, \texttt{xcol},\newline\texttt{row}
  & Referencia a una celda en la hoja ligada a \texttt{sheet} en
    \texttt{*sheet-env*}. \\
\midrule
\texttt{xl-expr-source-row}
  & \texttt{table-id},\newline\texttt{offset}
  & Fila de una tabla origen mapeada desde la fila l\'ogica actual.
    \texttt{offset} selecciona la subcelda dentro de una fila compuesta. \\
\midrule
\texttt{xl-expr-sheet-id}
  & \texttt{sheet}
  & Identificador can\'onico de la hoja ligada a \texttt{sheet} en
    \texttt{*sheet-env*}. \\
\midrule
\texttt{xl-expr-collect-over}
  & \texttt{groups},\newline\texttt{sheet-var},\newline\texttt{body}
  & Itera \texttt{body} sobre cada hoja de \texttt{groups}, ligando
    \texttt{sheet-var}, y concatena los resultados como lista separada
    por comas. \\
\bottomrule
\end{tabular}
\end{table}

% ............................................................................
\subsection{Formato condicional y posicionamiento}
\label{sec:ast-formato}
% ............................................................................

Estos nodos permiten expresar el coloreado condicional de celdas y la
colocaci\'on de expresiones en posiciones relativas a las tablas.  La
separaci\'on en nodos propios garantiza que el DSL no exponga coordenadas
absolutas: el backend resuelve las posiciones finales en funci\'on del
tama\~no real de los datos.

\begin{table}[h!]
\centering
\caption{Nodos de formato condicional y posicionamiento relativo.}
\label{tab:ast-formato}
\renewcommand{\arraystretch}{1.3}
\begin{tabular}{p{3.6cm} p{4.8cm} p{5.4cm}}
\toprule
\textbf{Nodo} & \textbf{Atributos} & \textbf{Papel} \\
\midrule
\texttt{xl-style-rule}
  & \texttt{rule-condition},\newline
    \texttt{target-columns}
  & Regla de coloreado condicional.  \texttt{rule-condition}
    es un nodo de expresi\'on booleana; \texttt{target-columns}
    enumera las columnas que se colorean cuando la condici\'on
    es verdadera. \\
\midrule
\texttt{xl-anchor}
  & \texttt{table-id},\newline
    \texttt{col-name},\newline
    \texttt{anchor-type}
  & Punto de referencia dentro de una tabla.
    \texttt{anchor-type} es \texttt{:ultima-fila} o
    \texttt{:primera-fila}. \\
\midrule
\texttt{xl-nav}
  & \texttt{anchor}, \texttt{steps}
  & Posici\'on calculada como un ancla m\'as una secuencia de
    pasos direccionales (arriba, abajo, izquierda, derecha). \\
\midrule
\texttt{xl-sheet-fixed-expr}
  & \texttt{pos}, \texttt{expr}
  & Expresi\'on colocada en una posici\'on relativa de la hoja.
    \texttt{pos} es un nodo \texttt{xl-nav}. \\
\bottomrule
\end{tabular}
\end{table}

% ---------------------------------------------------------------------------
\section{Definici\'on de tablas: \texttt{def-table}}
\label{sec:def-table}
% ---------------------------------------------------------------------------

La forma central del DSL para definir datos tabulares es \texttt{def-table}.
Produce una \emph{plantilla} de tabla reutilizable: una descripci\'on
abstracta de las columnas, sus cabeceras, las f\'ormulas que se calculan a
partir de ellas y la regla de coloreado que se aplica sobre sus filas.  La
plantilla no contiene datos; se instancia con datos concretos al ensamblar
el workbook (secci\'on~\ref{sec:workbook}).

\begin{lstlisting}[language=CommonLisp,
  caption={Sintaxis completa de def-table.}]
(def-table nombre-tabla
  ((col1 "Cabecera1") (col2 "Cabecera2") ...)  ; columnas
  :cell-height N       ; filas por fila logica (default 1)
  :cell-width  N       ; columnas por columna logica (default 1)
  :computed            ; formulas derivadas (opcional)
    ((colN (expresion-dsl)))
  :render              ; regla de coloreado (opcional)
    (conditional-rendering
      :condition (expresion-booleana)
      :target-columns (col1 col2 ...)))
\end{lstlisting}

Cada par \texttt{(col "Cabecera")} declara un nombre de columna DSL
(s\'{\i}mbolo) y el texto de su encabezado visual.  El orden de las columnas
es el orden f\'{\i}sico en el documento de salida.

La clave \texttt{:computed} asocia columnas a expresiones DSL que el
compilador traduce a f\'ormulas.  Los valores iniciales de esas columnas
en los datos de entrada se ignoran: la celda siempre mostrar\'a el resultado
de la f\'ormula calculada.

La clave \texttt{:render} acepta una forma \texttt{conditional-rendering}
que especifica bajo qu\'e condici\'on se colorean determinadas columnas.
El mecanismo concreto lo decide el backend.

El listado siguiente define la tabla de profesores del caso de uso
principal, con columna calculada y regla de coloreado condicional:

\begin{lstlisting}[language=CommonLisp,
  caption={Tabla de profesores con columna computed y regla de coloreado.}]
(def-table profesores-table
  ((nombre "") (grado "") (defensas ""))
  :cell-height 1
  :cell-width  1
  :computed
    ;; Cuenta cuantas veces aparece este profesor en los cinco roles
    ((defensas (countif (trange defensas-table tutor secretario)
                        (col nombre))))
  :render
    (conditional-rendering
      :condition
        (exists (r :rows-of defensas-table)
          :self-in  (tutor oponente presidente vocal secretario)
          :matching (dia hora))
      :target-columns (nombre)))
\end{lstlisting}

La variable \texttt{r} en la forma \texttt{exists} es un alias documental
de la fila iterada, an\'alogo a un alias de tabla en SQL.  La regla de
coloreado declara que la celda \texttt{nombre} debe resaltarse cuando el
profesor participa en alguna defensa con conflicto de horario; el mecanismo
lo decide el backend.

% ---------------------------------------------------------------------------
\section{Expresiones de celda}
\label{sec:expresiones-dsl}
% ---------------------------------------------------------------------------

El DSL ofrece formas simb\'olicas para construir expresiones de celda sin
escribir cadenas de f\'ormula directamente.  El compilador resuelve los
nombres de columna a sus coordenadas reales, los rangos a sus l\'{\i}mites
absolutos y las operaciones a su sintaxis en el formato de salida; el
programador solo declara la intenci\'on.

La tabla~\ref{tab:formas-expresion} recoge las formas disponibles.

\begin{table}[h!]
\centering
\caption{Formas de expresi\'on del DSL y su significado.}
\label{tab:formas-expresion}
\renewcommand{\arraystretch}{1.3}
\begin{tabular}{p{5.5cm} p{8.3cm}}
\toprule
\textbf{Forma DSL} & \textbf{Significado} \\
\midrule
\texttt{(col nombre)}
  & Valor de la columna \texttt{nombre} en la fila actual. \\
\midrule
\texttt{(trange tabla-id desde hasta)}
  & Rango que abarca desde la columna \texttt{desde} hasta la
    columna \texttt{hasta} de \texttt{tabla-id}, para todas
    las filas de datos. \\
\midrule
\texttt{(countif rango criterio)}
  & Conteo condicional sobre un rango. \\
\midrule
\texttt{(counta rango)}
  & N\'umero de celdas no vac\'{\i}as en el rango. \\
\midrule
\texttt{(sum-range rango)}
  & Suma de los valores del rango. \\
\midrule
\texttt{(str "texto")}
  & Literal de texto fijo. \\
\midrule
\texttt{(\_equals a b)}
  & Igualdad entre dos expresiones. \\
\midrule
\texttt{(\_and a b)}, \texttt{(\_or a b)}
  & Conjunci\'on y disyunci\'on. \\
\midrule
\texttt{(nav ancla dir n \ldots)}
  & Posici\'on relativa: ancla m\'as pasos direccionales. \\
\bottomrule
\end{tabular}
\end{table}

La forma \texttt{trange} produce \emph{referencias cross-table}: permite que
una expresi\'on en \texttt{profesores-table} cite columnas de
\texttt{defensas-table} sin que el programador conozca las coordenadas
reales.  El compilador las resuelve en funci\'on del orden de declaraci\'on
de las tablas en la regi\'on.

\begin{lstlisting}[language=CommonLisp,
  caption={Expresiones DSL y sus f\'ormulas generadas (backend Excel).}]
;; (col nombre) en profesores-table -> columna K, fila relativa
(col nombre)          ;; -> K4  (en la fila 4)

;; (trange defensas-table tutor secretario) -> rango absoluto de los cinco roles
(trange defensas-table tutor secretario)   ;; -> $B$4:$F$6

;; La formula COUNTIF resultante
(countif (trange defensas-table tutor secretario) (col nombre))
;; -> COUNTIF($B$4:$F$6, K4)
\end{lstlisting}

% ---------------------------------------------------------------------------
\section{El cuantificador existencial \texttt{exists}}
\label{sec:exists-dsl}
% ---------------------------------------------------------------------------

El cuantificador \texttt{exists} permite declarar condiciones de
\emph{existencia} sobre filas de una tabla.  Preguntas del tipo ``existe
otra defensa con el mismo horario'' o ``aparece este profesor en alguna
defensa con conflicto'' son muy frecuentes en problemas de planificaci\'on,
pero expresarlas mediante disyunciones y bucles expl\'{\i}citos produce
c\'odigo verboso y dif\'{\i}cil de leer.  Con \texttt{exists} el
programador declara la intenci\'on sem\'antica a alto nivel; el backend se
encarga de generar la comprobaci\'on correcta.

El cuantificador existe en dos variantes seg\'un el dominio de b\'usqueda.

% ............................................................................
\subsection{Variante same-table: \texttt{:all-rows}}
\label{sec:exists-all-rows}
% ............................................................................

Esta variante busca otra fila dentro de la \emph{misma} tabla que satisfaga
condiciones de coincidencia de valores y solapamiento en columnas de rol.
Se usa para detectar si una fila del propio conjunto tiene alg\'un conflicto
con otra fila del mismo conjunto; por ejemplo, dos defensas programadas en
el mismo d\'{\i}a y hora con al menos un integrante de tribunal en com\'un.

\begin{lstlisting}[language=CommonLisp,
  caption={Cuantificador exists sobre la tabla actual.}]
(exists (r :all-rows)
  :matching    (dia hora)              ; mismos valores en ambas filas
  :any-overlap (tutor oponente presidente vocal secretario))
                                       ; algun valor compartido en estos slots
\end{lstlisting}

La sem\'antica es: ``existe otra fila de esta tabla (distinta de la actual)
cuya columna \texttt{dia} y cuya columna \texttt{hora} coinciden con los de
la fila actual, y que adem\'as comparte al menos un valor en las columnas
de rol''.

En el backend Excel, esta expresi\'on se traduce a una f\'ormula SUMPRODUCT
que multiplica tres factores: coincidencia en las columnas de agrupaci\'on,
exclusi\'on de la fila actual y solapamiento en al menos una columna de rol.

% ............................................................................
\subsection{Variante cross-table: \texttt{:rows-of}}
\label{sec:exists-rows-of}
% ............................................................................

Esta variante busca en las filas de \emph{otra} tabla, relacionando el
valor de la celda actual con columnas de esa tabla.  Se usa cuando el
elemento evaluado (p.\,ej.\ un profesor) y los datos de planificaci\'on
(p.\,ej.\ las defensas) residen en tablas distintas.

\begin{lstlisting}[language=CommonLisp,
  caption={Cuantificador exists sobre otra tabla (cross-table).}]
(exists (r :rows-of defensas-table)
  :self-in  (tutor oponente presidente vocal secretario)
               ; el valor de la celda actual aparece en alguno de estos slots
  :matching (dia hora))
               ; ese slot tiene mas de una defensa programada
\end{lstlisting}

La sem\'antica es: ``existe alguna fila de \texttt{defensas-table} en la que
el valor de la celda actual (el nombre del profesor) aparece en una de las
cinco columnas de rol, y cuyo par (\texttt{dia}, \texttt{hora}) tiene m\'as
de una defensa programada''.

En ambas variantes el alias de fila es puramente documental: el backend
genera directamente una f\'ormula vectorial y no eval\'ua el s\'{\i}mbolo
como variable de programa.  Su papel es an\'alogo al de un alias de tabla
en una sentencia SQL.

% ---------------------------------------------------------------------------
\section{Formato condicional declarativo}
\label{sec:formato-condicional-dsl}
% ---------------------------------------------------------------------------

El coloreado de celdas es una herramienta habitual en los documentos de
planificaci\'on: resaltar en rojo los profesores con conflicto permite
detectar problemas de un vistazo.  El DSL ofrece la forma
\texttt{conditional-rendering} para declarar ese coloreado sin escribir
f\'ormulas ni llamadas a la API de formato.

\begin{lstlisting}[language=CommonLisp,
  caption={Sintaxis de conditional-rendering.}]
(conditional-rendering
  :condition      (expresion-booleana)
  :target-columns (col1 col2 ...))
\end{lstlisting}

Cuando el nodo condici\'on es una instancia de \texttt{xl-expr-exists}, el
backend del caso de demostraci\'on produce una \texttt{FormulaRule} de Excel
incrustada en el archivo \texttt{.xlsx}, en lugar de colorear celdas
est\'aticamente.  Esta regla persiste en el libro y se reevalua
autom\'aticamente cada vez que el usuario edita los datos, sin necesidad de
volver a ejecutar el DSL.

La decisi\'on de usar \texttt{FormulaRule} en lugar de colorear celdas
directamente es exclusiva del backend: el programador del DSL no indica
qu\'e mecanismo de formato se utilizar\'a.

% ---------------------------------------------------------------------------
\section{Ensamblaje del workbook}
\label{sec:workbook}
% ---------------------------------------------------------------------------

Una vez definidas las plantillas de tabla, el programador ensambla el
workbook con tres macros: \texttt{libro} crea el nodo ra\'{\i}z del AST
con el nombre del archivo de salida; \texttt{hoja} agrupa tablas y
expresiones fijas en una pesta\~na del libro; \texttt{tabla} instancia
una plantilla con datos concretos.

\begin{lstlisting}[language=CommonLisp,
  caption={Ensamblaje del workbook con libro, hoja y tabla.}]
(libro nombre-libro
  :filename "Archivo.xlsx"
  :hojas (list
    (hoja "NombreHoja"
      (tabla plantilla1 :data *DATOS-1*)
      (tabla plantilla2 :data *DATOS-2*)
      :fixed-expressions
        (((nav ...) (expresion))
         ((nav ...) (expresion))))))
\end{lstlisting}

Cuando una \texttt{hoja} recibe m\'as de una \texttt{tabla}, el generador
las coloca en la misma regi\'on horizontal, separadas por una columna en
blanco.  Este dise\~no permite que las referencias cross-table
(\texttt{trange}) se resuelvan con coordenadas absolutas correctas.

La clave \texttt{:fixed-expressions} acepta pares de
\texttt{(posici\'on, expresi\'on)}.  La posici\'on se especifica con
\texttt{nav}; su valor se recalcula seg\'un el tama\~no real de los datos,
de modo que los totales de pie de tabla se desplazan autom\'aticamente si se
a\~naden m\'as filas.

% ---------------------------------------------------------------------------
\section{Posicionamiento relativo: \texttt{nav}}
\label{sec:nav-dsl}
% ---------------------------------------------------------------------------

Las expresiones fijas de una hoja no pueden tener coordenadas absolutas
codificadas, pues esas coordenadas dependen de cu\'antas filas de datos
tenga cada tabla en cada ejecuci\'on.  La forma \texttt{nav} resuelve esto:
calcula la posici\'on de una celda como un \emph{ancla} fija en una tabla
m\'as una secuencia de pasos direccionales.

Las anclas disponibles son \texttt{ultima-fila} y \texttt{primera-fila}.
A partir del ancla se encadenan pasos \texttt{abajo}, \texttt{arriba},
\texttt{izquierda} y \texttt{derecha} con su cantidad:

\begin{lstlisting}[language=CommonLisp,
  caption={Sintaxis de nav y anclas de tabla.}]
;; Una fila debajo de la ultima fila de datos
(nav (ultima-fila tabla-id col-name) abajo 1)

;; Una fila debajo y dos columnas a la derecha
(nav (ultima-fila tabla-id col-name) abajo 1 derecha 2)

;; Una fila encima de la primera fila de datos
(nav (primera-fila tabla-id col-name) arriba 1)
\end{lstlisting}

El compilador resuelve el ancla a la coordenada correcta en tiempo de
generaci\'on y produce una referencia absoluta que siempre apunta al lugar
correcto, independientemente del n\'umero de filas de datos.

% ---------------------------------------------------------------------------
\section{Generaci\'on y ejecuci\'on}
\label{sec:generacion-dsl}
% ---------------------------------------------------------------------------

Las dos \'ultimas formas del DSL cierran el ciclo de compilaci\'on:
\texttt{xl-generate} recorre el \'arbol AST y produce el programa de
construcci\'on del libro; \texttt{xl-run-generated} ejecuta ese programa
y crea el archivo de salida.

\begin{lstlisting}[language=CommonLisp,
  caption={Compilaci\'on del AST y ejecuci\'on.}]
;; Recorre el AST y genera el programa de construccion del libro
(xl-generate horario-tesis "horario-tesis.py")

;; Ejecuta el programa generado y crea el archivo .xlsx
(xl-run-generated "horario-tesis.py")
\end{lstlisting}

% ---------------------------------------------------------------------------
\section{Ejemplo completo: Horario de Defensas de Tesis}
\label{sec:ejemplo-completo}
% ---------------------------------------------------------------------------

El caso de uso principal del DSL es la generaci\'on autom\'atica de un
horario de defensas de tesis con dos tablas relacionadas.  La primera tabla,
\texttt{defensas-table}, registra cada defensa con su tribunal (cinco roles:
tutor, oponente, presidente, vocal y secretario), su slot horario y el local.
La segunda tabla, \texttt{profesores-table}, enumera los profesores con su
grado acad\'emico y el n\'umero de defensas en que participan; sus filas se
colorean en rojo cuando el profesor tiene un conflicto de horario.

El listado siguiente es el programa DSL completo; en \'el se observan todas
las construcciones descritas en el cap\'{\i}tulo.

\begin{lstlisting}[language=CommonLisp,
  caption={Programa DSL completo: Horario de Defensas de Tesis.},
  label={lst:tutorial}]
(load "dsl-directo.lisp")
(load "data-tutorial.lisp")   ; *DATOS-DEFENSAS*, *DATOS-PROFESORES*

;; -- TABLA 1: defensas ------------------------------------------------
;; Nueve columnas: cinco de rol (tutor...secretario) mas dia, hora, local.
;; Las cinco columnas de rol son contiguas para poder referenciarlas como
;; rango con (trange defensas-table tutor secretario).
(def-table defensas-table
  ((estudiante "") (tutor "") (oponente "") (presidente "")
   (vocal "") (secretario "") (dia "") (hora "") (local ""))
  :cell-height 1
  :cell-width 1)

;; -- TABLA 2: profesores ----------------------------------------------
;; La columna "defensas" cuenta apariciones del nombre en los cinco roles.
;; El formato condicional marca en rojo a quienes tienen conflicto.
(def-table profesores-table
  ((nombre "") (grado "") (defensas ""))
  :cell-height 1
  :cell-width  1
  :computed
    ((defensas (countif (trange defensas-table tutor secretario)
                        (col nombre))))
  :render
    (conditional-rendering
      :condition
        (exists (r :rows-of defensas-table)
          :self-in  (tutor oponente presidente vocal secretario)
          :matching (dia hora))
      :target-columns (nombre)))

;; -- WORKBOOK ---------------------------------------------------------
(libro horario-tesis
  :filename "Horario_Tesis.xlsx"
  :hojas (list
    (hoja "Defensas"
      (tabla defensas-table   :data *DATOS-DEFENSAS*)
      (tabla profesores-table :data *DATOS-PROFESORES*)
      :fixed-expressions
        (((nav (ultima-fila defensas-table estudiante) abajo 1)
          (str "Total defensas:"))
         ((nav (ultima-fila defensas-table estudiante) abajo 1 derecha 1)
          (counta (trange defensas-table estudiante)))
         ((nav (ultima-fila profesores-table nombre) abajo 1)
          (str "Total profesores:"))
         ((nav (ultima-fila profesores-table nombre) abajo 1 derecha 2)
          (sum-range (trange profesores-table defensas)))))))

(xl-generate horario-tesis "horario-tesis.py")
(xl-run-generated "horario-tesis.py")
\end{lstlisting}

Los datos de entrada contienen tres defensas y seis profesores.  Dos de las
defensas comparten el slot \texttt{Lunes 10:00} con roles solapados.
El programa genera \texttt{Horario\_Tesis.xlsx} con la siguiente estructura:

\begin{itemize}
  \item Columnas de \texttt{defensas-table} con tres filas de datos.
  \item Columna separadora en blanco de la regi\'on.
  \item Columnas de \texttt{profesores-table} con seis filas de datos
    y la columna \texttt{Defensas} calculada mediante \texttt{COUNTIF}.
  \item Fila de totales debajo de cada tabla con posicionamiento din\'amico.
  \item Regla de formato condicional aplicada a la columna \texttt{nombre}:
    los profesores que participan en al menos una defensa cuyo slot tiene
    m\'as de una defensa programada aparecen resaltados en rojo.
\end{itemize}

La regla de coloreado implementa exactamente la sem\'antica del
\texttt{exists} declarado en el DSL.  Al ser una \texttt{FormulaRule}
incrustada en el archivo, cualquier edici\'on posterior de los datos provoca
una reevaluaci\'on autom\'atica del coloreado, sin necesidad de volver a
ejecutar el DSL.

% ===========================================================================
%  Capítulo 4 — Generación de código
% ===========================================================================
\chapter{Generaci\'on de c\'odigo}
\label{cap:generacion}

La generaci\'on de c\'odigo es la etapa del sistema que traduce el AST
construido por las formas del DSL a un programa Python que produce el archivo
\xlsx{}.  Este cap\'{\i}tulo describe c\'omo el backend resuelve las
coordenadas Excel, c\'omo compila cada tipo de nodo de expresi\'on a una
cadena de f\'ormula, y c\'omo genera las f\'ormulas SUMPRODUCT que
implementan el formato condicional persistente.

% ---------------------------------------------------------------------------
\section{Arquitectura del backend}
\label{sec:arquitectura-backend}
% ---------------------------------------------------------------------------

El backend est\'a organizado en dos archivos con responsabilidades separadas:

\begin{description}
  \item[\file{code-generation-utils.lisp}.] Funciones auxiliares, variables
    din\'amicas de configuraci\'on y la funci\'on de entrada principal.
    Incluye: \dsl{col->letter} (conversi\'on de n\'umero de columna a letra
    Excel), \dsl{build-col-map} (construcci\'on del mapa s\'{\i}mbolo~DSL
    $\to$ letra), \dsl{xl-generate} y \dsl{xl-run-generated}.

  \item[\file{generate-code-direct.lisp}.] Los m\'etodos CLOS de la
    funci\'on gen\'erica \dsl{compile-excel-formula} y los m\'etodos del
    generador de c\'odigo Python.  Este archivo es el \'unico lugar del
    sistema donde pueden aparecer letras de columna Excel, cadenas de
    f\'ormula y llamadas a openpyxl.
\end{description}

El principio de separaci\'on de responsabilidades queda codificado en la
propia arquitectura: \file{ast-def.lisp} declara \emph{qu\'e existe},
\file{dsl-directo.lisp} declara \emph{c\'omo se escribe} y
\file{generate-code-direct.lisp} declara \emph{a qu\'e se traduce}.

% ---------------------------------------------------------------------------
\section{Resoluci\'on de coordenadas}
\label{sec:resolucion-coordenadas}
% ---------------------------------------------------------------------------

% ............................................................................
\subsection{Conversi\'on de n\'umero de columna a letra Excel}
\label{ssec:col-letter}
% ............................................................................

La funci\'on \dsl{col->letter} convierte un entero $n$ (la posici\'on
ordinal de una columna, con base~1) a la letra o letras de columna
correspondientes en Excel: columna~1~$\to$~\texttt{A}, columna~26~$\to$~\texttt{Z},
columna~27~$\to$~\texttt{AA}, etc.

\begin{lstlisting}[language=CommonLisp,
  caption={Funci\'on col-\textgreater{}letter.}]
(defun col->letter (n)
  (with-output-to-string (s)
    (loop while (> n 0)
          do (multiple-value-bind (q r) (floor (1- n) 26)
               (setf n q)
               (write-char (code-char (+ 65 r)) s)))))
\end{lstlisting}

El algoritmo es una conversi\'on de base 10 a base 26 con el ajuste
habitual ($\lfloor (n-1)/26 \rfloor$, $((n-1) \bmod 26) + 65$) que mapea
los residuos a los caracteres ASCII de las letras may\'usculas.

% ............................................................................
\subsection{Construcci\'on del mapa de columnas}
\label{ssec:col-map}
% ............................................................................

Para cada tabla de la regi\'on se construye un \emph{mapa de columnas}:
una lista de asociaci\'on que relaciona cada s\'{\i}mbolo DSL de columna con
su letra Excel absoluta.  La funci\'on \dsl{build-col-map} construye el mapa
de una tabla sola, asumiendo que empieza en la columna~1.  La funci\'on
\dsl{build-table-col-maps} construye el mapa global de toda la regi\'on,
teniendo en cuenta el desplazamiento de cada tabla dentro de la regi\'on:

\begin{lstlisting}[language=CommonLisp,
  caption={Construcci\'on del mapa de columnas de una regi\'on.}]
(defun build-col-map (tbl)
  (loop for name in (col-names tbl)
        for i from 1
        collect (cons name (col->letter i))))

(defun build-table-col-maps (tables offsets)
  (loop for tbl in tables
        for offset in offsets
        collect (cons (id tbl)
                      (loop for name in (col-names tbl)
                            for i from 1
                            collect (cons name
                                          (col->letter (+ offset i)))))))
\end{lstlisting}

El resultado se almacena en la variable din\'amica \dsl{*table-col-maps*},
disponible durante toda la compilaci\'on de la regi\'on.

% ............................................................................
\subsection{Cálculo de rangos de fila}
\label{ssec:rangos-fila}
% ............................................................................

Cada tabla tiene una primera fila de datos (\dsl{first-row}, configurable,
por defecto~4) y una \'ultima fila que depende del n\'umero de filas de datos
y de \dsl{cell-height}:

\begin{equation}
\textit{last-row} = \textit{first-row} + |\textit{datos}| \times \textit{cell-height} - 1
\end{equation}

Estos valores se almacenan en la variable din\'amica \dsl{*table-data-rows*}
como una lista de asociaci\'on \dsl{(table-id . (first-row . last-row))},
lo que permite que los nodos \dsl{xl-table-range} usen los rangos de fila
correctos para cada tabla aunque la compilaci\'on se realice dentro del
contexto de una tabla diferente.

% ---------------------------------------------------------------------------
\section{La funci\'on gen\'erica \texttt{compile-excel-formula}}
\label{sec:compile-excel-formula}
% ---------------------------------------------------------------------------

El m\'etodo central del backend es la funci\'on gen\'erica
\dsl{compile-excel-formula}, que toma un nodo AST y devuelve la cadena de
f\'ormula Excel correspondiente.  Su firma es:

\begin{lstlisting}[language=CommonLisp,
  caption={Declaraci\'on de compile-excel-formula.}]
(defgeneric compile-excel-formula
    (expr col-map data-names row-num first-row last-row)
  (:documentation
   "Compila un nodo AST de expresion a string de formula Excel"))
\end{lstlisting}

Los par\'ametros son:
\begin{description}
  \item[\dsl{expr}] El nodo AST a compilar.
  \item[\dsl{col-map}] El mapa de columnas de la tabla actual
    (\dsl{sym $\to$ letra}).
  \item[\dsl{data-names}] Lista de s\'{\i}mbolos DSL de la tabla actual
    (para \dsl{VLOOKUP}).
  \item[\dsl{row-num}] N\'umero de fila Excel que se est\'a compilando
    (para referencias relativas).
  \item[\dsl{first-row}] Primera fila de datos de la tabla actual.
  \item[\dsl{last-row}] \'Ultima fila de datos de la tabla actual.
\end{description}

El despacho m\'ultiple de CLOS garantiza que, para cada tipo de nodo, se
ejecuta el m\'etodo correcto sin ning\'un \dsl{cond} centralizado.
A\~nadir un nuevo tipo de expresi\'on al lenguaje s\'olo requiere:
\begin{enumerate}
  \item Definir una nueva clase con \dsl{defclass*} en \file{ast-def.lisp}.
  \item Definir un macro en \file{dsl-directo.lisp} que construya el nodo.
  \item Definir un \dsl{defmethod} de \dsl{compile-excel-formula} en
    \file{generate-code-direct.lisp}.
\end{enumerate}

Ning\'un c\'odigo existente necesita modificarse.

% ............................................................................
\subsection{M\'etodos para expresiones simples}
\label{ssec:metodos-simples}
% ............................................................................

\begin{lstlisting}[language=CommonLisp,
  caption={M\'etodos compile-excel-formula para expresiones simples.}]
;; Referencia a columna de la tabla actual: "K4"
(defmethod compile-excel-formula
    ((e clase-xl-expr-column-ref) col-map data-names
     row-num first-row last-row)
  (let* ((letter (cdr (assoc (name e) col-map)))
         (ctx (context e)))
    (if (typep ctx 'clase-xl-expr-previous-row)
        (format nil "~a~a" letter (1- row-num))
        (format nil "~a~a" letter row-num))))

;; Comparacion de igualdad: "K4=B4"
(defmethod compile-excel-formula
    ((e clase-xl-expr-equals) col-map data-names
     row-num first-row last-row)
  (format nil "~a=~a"
          (compile-excel-formula (a e) col-map data-names
                                  row-num first-row last-row)
          (compile-excel-formula (b e) col-map data-names
                                  row-num first-row last-row)))

;; Literal de texto: "\"Total\""
(defmethod compile-excel-formula
    ((e clase-xl-expr-string) col-map data-names
     row-num first-row last-row)
  (format nil "~s" (value e)))
\end{lstlisting}

% ............................................................................
\subsection{M\'etodos para rangos y agregados}
\label{ssec:metodos-rangos}
% ............................................................................

\begin{lstlisting}[language=CommonLisp,
  caption={M\'etodo para xl-table-range.}]
;; Rango absoluto de columnas de una tabla especifica de la region.
;; Usa *table-data-rows* para obtener los limites exactos de esa tabla.
(defmethod compile-excel-formula
    ((e clase-xl-table-range) col-map data-names
     row-num first-row last-row)
  (let* ((tmap        (cdr (assoc (table-id e) *table-col-maps*)))
         (from-letter (cdr (assoc (name (from-col e)) tmap)))
         (to-letter   (if (to-col e)
                          (cdr (assoc (name (to-col e)) tmap))
                          from-letter))
         (tbl-rows    (cdr (assoc (table-id e) *table-data-rows*)))
         (act-first   (if tbl-rows (car tbl-rows) first-row))
         (act-last    (if tbl-rows (cdr tbl-rows) last-row)))
    (format nil "$~a$~a:$~a$~a"
            from-letter act-first to-letter act-last)))
\end{lstlisting}

La clave t\'ecnica de este m\'etodo es que consulta \dsl{*table-data-rows*}
con el identificador de la \emph{tabla referenciada}, no con el de la tabla
actual.  As\'{\i}, \dsl{(trange defensas-table tutor secretario)} compila al
rango de \dsl{defensas-table} aunque se est\'e compilando una f\'ormula de
\dsl{profesores-table}.

% ---------------------------------------------------------------------------
\section{Generaci\'on de formato condicional}
\label{sec:cf-generacion}
% ---------------------------------------------------------------------------

El formato condicional es el aspecto t\'ecnicamente m\'as elaborado de la
generaci\'on de c\'odigo.  El backend recibe un nodo \dsl{xl-expr-exists} y
produce una f\'ormula SUMPRODUCT que implementa la sem\'antica del
cuantificador existencial declarado.

% ............................................................................
\subsection{Variante same-table: SUMPRODUCT con exclusi\'on de fila actual}
\label{ssec:cf-same-table}
% ............................................................................

Cuando \dsl{exists} usa \dsl{:all-rows}, el backend genera una f\'ormula
SUMPRODUCT que, para cada fila, multiplica tres factores:
\begin{enumerate}
  \item Coincidencia en las columnas de agrupaci\'on (\dsl{:matching}).
  \item Exclusi\'on de la fila actual (para no comparar una fila consigo misma).
  \item Solapamiento en al menos una columna de rol (\dsl{:any-overlap}).
\end{enumerate}

Para una regla con \dsl{:matching (dia hora)} y
\dsl{:any-overlap (tutor \ldots\ secretario)}, la f\'ormula generada para
la fila~4 es:

\begin{lstlisting}[style=pythonstyle,
  caption={F\'ormula Excel CF para la variante same-table (fila 4).}]
=SUMPRODUCT(
    ($G$4:$G$6=$G4) *          -- dia coincide en ambas filas
    ($H$4:$H$6=$H4) *          -- hora coincide en ambas filas
    (ROW($B$4:$B$6)<>ROW($B4)) *  -- excluir la fila actual
    (($B$4:$B$6=$B4)+($C$4:$C$6=$C4)+...+($F$4:$F$6=$F4))
) > 0
\end{lstlisting}

Las referencias con columna fija y fila relativa (p.\,ej.\ \texttt{\$G4})
hacen que la f\'ormula se eval\'ue correctamente para cada fila del rango:
al desplazarse de la fila~4 a la fila~5, las referencias de ``fila actual''
avanzan autom\'aticamente.  El resultado de \texttt{SUMPRODUCT} es el n\'umero
de filas del dominio que satisfacen las tres condiciones; la comparaci\'on
\texttt{> 0} convierte ese conteo en un valor booleano para la regla de
formato condicional.

% ............................................................................
\subsection{Variante cross-table: SUMPRODUCT con COUNTIFS}
\label{ssec:cf-cross-table}
% ............................................................................

Cuando \dsl{exists} usa \dsl{:rows-of}, el backend genera una f\'ormula
SUMPRODUCT que combina dos condiciones:
\begin{enumerate}
  \item Que el valor de la celda actual aparezca en al menos una columna de
    rol de la tabla referenciada (\dsl{:self-in}).
  \item Que el slot de tiempo de esa fila tenga m\'as de una defensa
    programada (\dsl{:matching} con \texttt{COUNTIFS}).
\end{enumerate}

Para la regla de \dsl{profesores-table} con
\dsl{:self-in (tutor \ldots\ secretario)} y \dsl{:matching (dia hora)},
la f\'ormula generada para la celda K4 es:

\begin{lstlisting}[style=pythonstyle,
  caption={F\'ormula Excel CF para la variante cross-table (celda K4).}]
=SUMPRODUCT(
    (($B$4:$B$6=$K4)+($C$4:$C$6=$K4)+($D$4:$D$6=$K4)+
     ($E$4:$E$6=$K4)+($F$4:$F$6=$K4)) *
    (COUNTIFS($G$4:$G$6,$G$4:$G$6,
              $H$4:$H$6,$H$4:$H$6) > 1)
) > 0
\end{lstlisting}

La referencia \texttt{\$K4} corresponde al nombre del profesor en la celda
actual (columna~K bloqueada, fila relativa).  Los rangos
\texttt{\$B\$4:\$B\$6} a \texttt{\$F\$4:\$F\$6} son las columnas de rol de
\dsl{defensas-table}, completamente absolutas.

\texttt{COUNTIFS} con cada par de columnas de coincidencia iguales a s\'{\i}
mismas produce el n\'umero de filas que comparten exactamente los mismos
valores en esas columnas; comparar el resultado con~$1$ selecciona los
slots de tiempo con m\'as de una defensa.  El producto SUMPRODUCT de ambas
condiciones devuelve el n\'umero de filas de \dsl{defensas-table} en que el
profesor actual participa en un slot conflictivo.

Al ser una \texttt{FormulaRule} incrustada en el archivo Excel, cualquier
edici\'on posterior de los datos provoca una reevaluaci\'on autom\'atica del
coloreado, sin necesidad de volver a ejecutar el DSL.

% ---------------------------------------------------------------------------
\section{Generaci\'on de referencias cruzadas entre hojas}
\label{sec:cf-cross-sheet}
% ---------------------------------------------------------------------------

El nodo \dsl{xl-expr-collect-over} itera sobre una lista de hojas del
workbook y agrega la expresi\'on del cuerpo para cada una.  El mecanismo de
iteraci\'on usa la variable din\'amica \dsl{*sheet-env*}: durante la
compilaci\'on de cada t\'ermino, la variable de iteraci\'on (\dsl{sheet-var})
se liga temporalmente al nombre de la hoja actual.

\begin{lstlisting}[language=CommonLisp,
  caption={M\'etodo compile-excel-formula para collect-over.}]
(defmethod compile-excel-formula
    ((e clase-xl-expr-collect-over) col-map data-names
     row-num first-row last-row)
  (let* ((groups (groups e))
         (sh-var (sheet-var e))
         (body   (body e))
         (terms  (mapcar
                   (lambda (g)
                     ;; Liga sh-var -> g solo durante este termino
                     (let ((*sheet-env* (acons sh-var g *sheet-env*)))
                       (compile-excel-formula body col-map data-names
                                              row-num first-row last-row)))
                   groups))
         (inner  (format nil "~{~a~^&~}" terms)))
    (format nil "SUBSTITUTE(TRIM(~a),\" \",\",\")" inner)))
\end{lstlisting}

La f\'ormula resultante concatena los $N$ t\'erminos con \texttt{\&} y los
envuelve en \texttt{SUBSTITUTE(TRIM(\ldots))} para eliminar espacios extras
y reemplazarlos con comas, produciendo una lista de grupos separada por comas.

Para el horario de la facultad con 17 grupos, el m\'etodo genera una f\'ormula
que agrega la ocupaci\'on del Aula~1 en el turno~3 del lunes a trav\'es de
las 17 hojas del workbook, sin que el programador haya escrito ning\'un
nombre de hoja ni letra de columna.

% ---------------------------------------------------------------------------
\section{Generaci\'on del programa Python}
\label{sec:generacion-python}
% ---------------------------------------------------------------------------

El generador principal recorre el \'arbol del workbook y emite un programa
Python que usa \texttt{openpyxl} para construir el archivo \xlsx{}.  El
programa generado:

\begin{enumerate}
  \item Crea un workbook openpyxl.
  \item Para cada hoja, crea una hoja de trabajo y la configura.
  \item Para cada tabla de cada regi\'on, escribe las cabeceras y los datos.
  \item Para cada columna calculada, escribe la f\'ormula en cada fila de
    datos (compilada por \dsl{compile-excel-formula} en tiempo de generaci\'on
    del Python).
  \item Para cada regla de estilo condicional, crea una
    \texttt{FormulaRule} con la f\'ormula SUMPRODUCT generada y la a\~nade
    al rango correspondiente.
  \item Para cada expresi\'on fija, resuelve la posici\'on \dsl{nav} y escribe
    la expresi\'on compilada en la celda correspondiente.
  \item Guarda el workbook en el archivo \xlsx{} configurado.
\end{enumerate}

El listado siguiente muestra el fragmento del programa Python generado para
la columna calculada \dsl{defensas} de \dsl{profesores-table}:

\begin{lstlisting}[style=pythonstyle,
  caption={Fragmento del programa Python generado para la columna computed.}]
# Columna: defensas (profesores-table)
ws["M4"] = "=COUNTIF($B$4:$F$6,K4)"
ws["M5"] = "=COUNTIF($B$4:$F$6,K5)"
ws["M6"] = "=COUNTIF($B$4:$F$6,K6)"
ws["M7"] = "=COUNTIF($B$4:$F$6,K7)"
ws["M8"] = "=COUNTIF($B$4:$F$6,K8)"
ws["M9"] = "=COUNTIF($B$4:$F$6,K9)"

# FormulaRule para profesores-table, columna nombre
from openpyxl.formatting.rule import FormulaRule
from openpyxl.styles import PatternFill
red_fill = PatternFill(start_color="FF6666",
                       end_color="FF6666", fill_type="solid")
formula = ("SUMPRODUCT((($B$4:$B$6=$K4)+($C$4:$C$6=$K4)+"
           "($D$4:$D$6=$K4)+($E$4:$E$6=$K4)+($F$4:$F$6=$K4))*"
           "(COUNTIFS($G$4:$G$6,$G$4:$G$6,$H$4:$H$6,$H$4:$H$6)>1))>0")
ws.conditional_formatting.add(
    "K4:K9",
    FormulaRule(formula=[formula], fill=red_fill))
\end{lstlisting}

Las f\'ormulas de las celdas est\'an calculadas en tiempo de generaci\'on del
Python: el programa Python que se ejecuta s\'olo asigna cadenas de texto a las
celdas.  Las referencias relativas dentro de las f\'ormulas las eval\'ua Excel
al abrir el archivo.  La \texttt{FormulaRule}, en cambio, se reevalua
en Excel din\'amicamente cada vez que cambian los datos.

% ===========================================================================
%  Capítulo 5 — Detalles de implementación
% ===========================================================================
\chapter{Detalles de implementaci\'on}
\label{cap:implementacion}

Este cap\'{\i}tulo describe los aspectos internos del sistema: la organizaci\'on
de los archivos con sus responsabilidades, los macros que implementan el DSL,
el patr\'on de extensibilidad del backend y la validaci\'on del sistema sobre
los tres casos de uso presentados en el cap\'{\i}tulo~\ref{cap:dominio}.

% ---------------------------------------------------------------------------
\section{Organizaci\'on en cuatro archivos}
\label{sec:organizacion}
% ---------------------------------------------------------------------------

El sistema est\'a organizado en cuatro archivos Lisp con responsabilidades
estrictamente separadas.  La dependencia entre ellos es un\'{\i}voca:
cada archivo carga \'unicamente al que le precede.

\begin{table}[h!]
\centering
\caption{Archivos del sistema y sus responsabilidades.}
\label{tab:archivos}
\renewcommand{\arraystretch}{1.3}
\begin{tabular}{p{4.8cm} p{9.0cm}}
\toprule
\textbf{Archivo} & \textbf{Responsabilidad} \\
\midrule
\file{codigo-tesis.lisp}
  & Utilidades generales: \dsl{symb}, \dsl{mkstr}, \dsl{keywd},
    \dsl{ppme} y la macro \dsl{defclass*}.  No conoce ning\'un concepto
    del dominio de hojas de c\'alculo. \\
\midrule
\file{ast-def.lisp}
  & Definici\'on de todas las clases del AST con \dsl{defclass*}.
    Backend-agn\'ostico: no contiene ninguna letra de columna Excel,
    cadena de f\'ormula ni referencia a ninguna biblioteca de salida. \\
\midrule
\file{code-generation-utils.lisp} \newline
\file{generate-code-direct.lisp}
  & Backend Excel/Python.  \'Unico lugar del sistema donde viven letras
    de columna, cadenas de f\'ormula, estructuras openpyxl y conocimiento
    de la disposici\'on concreta del workbook.  El primero contiene
    las utilidades y variables din\'amicas; el segundo, los m\'etodos
    \dsl{compile-excel-formula}. \\
\midrule
\file{dsl-directo.lisp}
  & Macros del DSL: \dsl{def-table}, \dsl{tabla}, \dsl{hoja}, \dsl{hoja-v},
    \dsl{libro}, \dsl{exists}, \dsl{conditional-rendering}, \dsl{nav},
    \dsl{col}, \dsl{tcol}, \dsl{trange}, \dsl{countif}, \dsl{counta},
    \dsl{sum-range}, \dsl{str}, \dsl{concat}, \dsl{time-add}, y dem\'as.
    No contiene l\'ogica de generaci\'on de c\'odigo. \\
\bottomrule
\end{tabular}
\end{table}

Esta separaci\'on tiene consecuencias pr\'acticas importantes: es posible
cargar \file{ast-def.lisp} y construir un AST sin haber cargado el backend,
lo que facilita las pruebas de la sintaxis y la sem\'antica del DSL de forma
independiente.

% ---------------------------------------------------------------------------
\section{Macros del DSL}
\label{sec:macros-dsl}
% ---------------------------------------------------------------------------

Los macros del DSL se clasifican en tres grupos seg\'un su papel en la
jerarqu\'{\i}a del AST.

% ............................................................................
\subsection{Macros de expresi\'on}
\label{ssec:macros-expresion}
% ............................................................................

Los macros de expresi\'on construyen nodos de la familia de expresiones de
celda.  Todos son equivalencias directas: reciben los argumentos del
programador y los pasan al constructor CLOS correspondiente, haciendo el
quoting adecuado para los s\'{\i}mbolos del dominio.

\begin{lstlisting}[language=CommonLisp,
  caption={Macros de expresi\'on representativos.}]
;; (col nombre) -> xl-expr-column-ref con nombre quoteado
(defmacro col (name &optional context)
  (if context
      `(xl-expr-column-ref :name ',name :context ,context)
      `(xl-expr-column-ref :name ',name :context nil)))

;; (trange tabla desde) / (trange tabla desde hasta)
;; -> xl-table-range con los dos extremos como xl-expr-table-col-ref
(defmacro trange (table-id from &optional to)
  `(xl-table-range
     :table-id ',table-id
     :from-col (xl-expr-table-col-ref :table-id ',table-id
                                       :name ',from :context nil)
     :to-col   ,(if to
                    `(xl-expr-table-col-ref :table-id ',table-id
                                             :name ',to :context nil)
                    nil)))

;; (exists (var :rows-of tabla) :self-in (...) :matching (...))
;; -> xl-expr-exists con domain cross-table
(defmacro exists ((var &rest domain-spec)
                  &key matching any-overlap self-in)
  (let ((domain-type (first domain-spec)))
    (ecase domain-type
      (:all-rows
       `(xl-expr-exists :bind-var     ',var
                         :domain       (xl-domain-rows :table-id nil)
                         :match-keys   ',matching
                         :overlap-cols ',any-overlap
                         :self-in-cols  nil))
      (:rows-of
       (let ((table-id (second domain-spec)))
         `(xl-expr-exists :bind-var     ',var
                           :domain       (xl-domain-rows
                                           :table-id ',table-id)
                           :match-keys   ',matching
                           :overlap-cols ',any-overlap
                           :self-in-cols  ',self-in))))))
\end{lstlisting}

La regla de quoting de los macros es determinista: un s\'{\i}mbolo que es
nombre de columna o tabla en el dominio se quotea (\dsl{',sym}); un
s\'{\i}mbolo que es una variable Lisp en scope no se quotea (\dsl{,sym}).
Esta distinci\'on es expl\'{\i}cita en todos los macros.

% ............................................................................
\subsection{Macros de estructura}
\label{ssec:macros-estructura}
% ............................................................................

Los macros \dsl{def-table}, \dsl{hoja}, \dsl{hoja-v} y \dsl{libro}
construyen los nodos de la jerarqu\'{\i}a del workbook.

\dsl{def-table} es el m\'as complejo: genera una funci\'on que construye
un \dsl{xl-table} con sus slots rellenos.  Los \dsl{:inst-params} declarados
se a\~naden a la firma de esa funci\'on, lo que permite instanciar la misma
plantilla de tabla con valores de par\'ametro distintos:

\begin{lstlisting}[language=CommonLisp,
  caption={Expansi\'on simplificada de def-table.}]
(def-table aulas-dia-table
  ((dia "") (aula1 "") ...)
  :inst-params (dia)
  :computed ((aula1 (collect-over *grupos-all* (g) ...))))

;; Genera la funcion:
(defun aulas-dia-table (&key data params dia)
  (xl-table :id 'aulas-dia-table
            :col-names '(dia aula1 ...)
            :computed (list (cons 'aula1 (collect-over ...)))
            ...))

;; Que luego se instancia con:
(tabla aulas-dia-table :dia lun :data *datos-aulas-lunes*)
;; -> (aulas-dia-table :data *datos-aulas-lunes* :dia 'lun)
\end{lstlisting}

% ............................................................................
\subsection{El macro \texttt{nav} y el sistema de posicionamiento}
\label{ssec:macro-nav}
% ............................................................................

El macro \dsl{nav} convierte una secuencia plana de pares direcci\'on--cantidad
en una lista de cons-celdas:

\begin{lstlisting}[language=CommonLisp,
  caption={Implementaci\'on del macro nav.}]
(defmacro nav (anchor &rest steps)
  (let ((step-pairs
         (loop for (dir n) on steps by #'cddr
               collect `(cons ',dir ,n))))
    `(xl-nav :anchor ,anchor
             :steps (list ,@step-pairs))))

;; Ejemplo: (nav (ultima-fila t c) abajo 1 derecha 2)
;; Expande a:
;; (xl-nav :anchor (xl-anchor :table-id 't :col-name 'c
;;                              :anchor-type :ultima-fila)
;;          :steps (list (cons 'abajo 1) (cons 'derecha 2)))
\end{lstlisting}

El backend resuelve los pasos en tiempo de generaci\'on: dado el n\'umero de
fila de la \'ultima celda de datos, suma o resta las cantidades seg\'un la
direcci\'on y produce la coordenada Excel absoluta.

% ---------------------------------------------------------------------------
\section{Extensibilidad del sistema}
\label{sec:extensibilidad}
% ---------------------------------------------------------------------------

La extensibilidad es una consecuencia directa del uso de CLOS y de la
separaci\'on en archivos.  Considrese el ejemplo de a\~nadir soporte para
la funci\'on \texttt{AVERAGEIF} de Excel, que actualmente no est\'a en el
lenguaje.

\textbf{Paso 1 --- Nuevo nodo AST} (en \file{ast-def.lisp}):
\begin{lstlisting}[language=CommonLisp]
(defclass* xl-expr-averageif () (average-range criteria))
\end{lstlisting}

\textbf{Paso 2 --- Nuevo macro} (en \file{dsl-directo.lisp}):
\begin{lstlisting}[language=CommonLisp]
(defmacro averageif (range-expr criteria)
  `(xl-expr-averageif :average-range ,range-expr
                       :criteria ,criteria))
\end{lstlisting}

\textbf{Paso 3 --- Nuevo m\'etodo} (en \file{generate-code-direct.lisp}):
\begin{lstlisting}[language=CommonLisp]
(defmethod compile-excel-formula
    ((e clase-xl-expr-averageif) col-map data-names
     row-num first-row last-row)
  (let ((range-str (compile-excel-formula
                     (average-range e) col-map data-names
                     row-num first-row last-row))
        (criteria  (compile-excel-formula
                     (criteria e) col-map data-names
                     row-num first-row last-row)))
    (format nil "AVERAGEIF(~a,~a)" range-str criteria)))
\end{lstlisting}

Ning\'un c\'odigo existente necesita modificarse.  El macro \dsl{defclass*}
garantiza que la nueva clase tiene exactamente la misma interfaz que las dem\'as.

% ---------------------------------------------------------------------------
\section{Validaci\'on sobre los casos de uso}
\label{sec:validacion}
% ---------------------------------------------------------------------------

Los tres casos de uso descritos en el cap\'{\i}tulo~\ref{cap:dominio} se
validaron ejecutando el DSL y verificando manualmente el archivo
\xlsx{} generado.

% ............................................................................
\subsection{Horario de Defensas de Tesis}
\label{ssec:val-defensas}
% ............................................................................

El programa DSL del listado~\ref{lst:tutorial} genera \texttt{Horario\_Tesis.xlsx}
con la estructura esperada.  La verificaci\'on confirm\'o:

\begin{itemize}
  \item La columna \texttt{Defensas} de \dsl{profesores-table} muestra los
    valores correctos: Piad~=~2, Garc\'{\i}a~=~2, L\'opez~=~3, Torres~=~3,
    Ruiz~=~2, Soto~=~2.
  \item Los profesores L\'opez y Torres, que participan en las dos defensas
    del slot \texttt{Lunes~10:00}, aparecen resaltados en rojo.
  \item Los totales en la fila inferior a las tablas son correctos:
    ``Total defensas: 3'' y ``Total profesores: 6'' con las sumas esperadas.
  \item La \texttt{FormulaRule} persiste en el archivo: al cambiar
    manualmente la hora de la tercera defensa de \texttt{10:00} a
    \texttt{14:00}, el resaltado de L\'opez y Torres desaparece
    autom\'aticamente en Excel.
\end{itemize}

% ............................................................................
\subsection{Horario Docente de la Facultad}
\label{ssec:val-facultad}
% ............................................................................

El programa \file{ast-facultad.lisp} genera \texttt{Horario\_Facultad.xlsx}
con diecisiete hojas de grupo (Ciencia de la Computaci\'on en cuatro a\~nos,
Matem\'atica en cuatro a\~nos y Ciencia de Datos en cuatro a\~nos) m\'as una
hoja de ocupaci\'on de aulas.

La verificaci\'on confirm\'o:

\begin{itemize}
  \item Las columnas \dsl{asignadas} y \dsl{faltan} de \dsl{stats-table}
    calculan correctamente para cada grupo a partir del contenido de
    \dsl{turno-table}.
  \item La hoja de aulas muestra en cada celda los grupos que ocupan cada
    aula en cada turno, con la lista de grupos separada por comas producida
    por las f\'ormulas \texttt{SUBSTITUTE(TRIM(\ldots))}.
  \item Las celdas compuestas de \dsl{turno-table} (\dsl{cell-height~=~2})
    se generan correctamente: la subfila superior muestra la asignatura y
    la inferior el aula.
\end{itemize}

Este caso ejercita, adem\'as, la caracter\'{\i}stica de \dsl{inst-params}:
la tabla \dsl{aulas-dia-table} se define una sola vez y se instancia una
vez por d\'{\i}a de la semana, con \dsl{:dia lun}, \dsl{:dia mar}, etc.

% ............................................................................
\subsection{Discusi\'on comparativa}
\label{ssec:val-comparacion}
% ............................................................................

Para ilustrar el beneficio del DSL, se estima la cantidad de c\'odigo
necesaria para implementar el caso de defensas de forma imperativa directa
en Python con openpyxl, en comparaci\'on con el programa DSL.

La versi\'on imperativa directa requiere calcular manualmente las letras de
columna, construir las cadenas de f\'ormula COUNTIF con interpolaci\'on de
texto, y ensamblar la f\'ormula SUMPRODUCT con referencias absolutas calculadas
a partir del n\'umero de filas.  Una implementaci\'on representativa ocupa
aproximadamente 80--100 l\'{\i}neas de Python.

El programa DSL equivalente (listado~\ref{lst:tutorial}) ocupa 40 l\'{\i}neas,
de las cuales la mayor parte son la declaraci\'on de datos
(\dsl{def-table}, \dsl{libro}) y la expresi\'on de intenci\'on
(\dsl{exists}, \dsl{conditional-rendering}).  Ning\'una l\'{\i}nea
contiene una letra de columna, un \'{\i}ndice de fila o una cadena de f\'ormula.

M\'as importante que el n\'umero de l\'{\i}neas es que el programa DSL es
resistente a cambios estructurales: a\~nadir una columna a \dsl{defensas-table}
s\'olo requiere a\~nadir un par \dsl{(nombre "Cabecera")} a la lista de columnas;
el compilador recalcula autom\'aticamente todas las letras y rangos.  En la
versi\'on imperativa, el mismo cambio requiere actualizar manualmente todas
las referencias a columnas afectadas.

% ===========================================================================
%  Conclusiones y Recomendaciones
% ===========================================================================
\chapter*{Conclusiones}
\addcontentsline{toc}{chapter}{Conclusiones}
\markboth{Conclusiones}{Conclusiones}

El trabajo alcanz\'o los objetivos propuestos.  Se dise\~n\'o e implement\'o
un lenguaje de dominio espec\'{\i}fico en Common Lisp para la descripci\'on
declarativa de estructuras tabulares de planificaci\'on, y se valid\'o sobre
tres casos reales de la Facultad de Matem\'atica y Computaci\'on de la
Universidad de La Habana mediante un backend que produce archivos \xlsx{}.

Las conclusiones principales son las siguientes.

\begin{enumerate}

  \item \textbf{La separaci\'on entre descripci\'on sem\'antica y backend de
    generaci\'on es alcanzable con un DSL interno en Lisp.}  El sistema
    permite declarar tablas relacionadas, columnas calculadas, cuantificadores
    existenciales y reglas de coloreado condicional sin hacer referencia a
    ning\'un detalle concreto del sistema de salida: ning\'un \'{\i}ndice de
    columna, ninguna cadena de f\'ormula, ninguna llamada a API alguna.  El
    backend es el \'unico punto del sistema donde viven esos detalles, y es
    intercambiable sin modificar el programa DSL.

  \item \textbf{Las macros de Common Lisp son el mecanismo adecuado para
    implementar el DSL.}  La homoiconicidad del lenguaje hace que los macros
    sean transformadores de c\'odigo completamente naturales: cada forma del
    DSL es una lista anidada que el macro convierte en constructores CLOS,
    sin ning\'un parser externo.

  \item \textbf{El despacho m\'ultiple de CLOS ofrece extensibilidad sin
    modificar el c\'odigo existente.}  A\~nadir un nuevo tipo de expresi\'on
    al lenguaje requiere \'unicamente tres acciones: definir la clase del nodo
    con \dsl{defclass*}, definir el macro de superficie y definir el m\'etodo
    del backend.  El patr\'on se replic\'o consistentemente en los m\'as de
    veinte tipos de nodo implementados.

  \item \textbf{El cuantificador \dsl{exists} captura correctamente las
    condiciones de conflicto de planificaci\'on.}  Las dos variantes
    (\dsl{:all-rows} y \dsl{:rows-of}) cubren los patrones de conflicto
    m\'as frecuentes en los documentos analizados.  El backend genera
    f\'ormulas SUMPRODUCT que implementan esas sem\'anticas de forma robusta
    y producen \texttt{FormulaRule} persistentes en el \xlsx{}.

  \item \textbf{El posicionamiento relativo con \dsl{nav} elimina las
    coordenadas absolutas de los programas del DSL.}  Totales, etiquetas y
    celdas de resumen se colocan siempre en la posici\'on correcta
    independientemente del n\'umero de filas de datos, lo que hace al
    programa resiliente ante cambios en el tama\~no de los conjuntos.

  \item \textbf{El backend implementado genera archivos \xlsx{} con formato
    condicional activo.}  La decisi\'on de usar \texttt{FormulaRule} en lugar
    de coloreado est\'atico demuestra que la arquitectura es capaz de producir
    salidas con comportamiento din\'amico: la hoja de c\'alculo sigue siendo
    funcional cuando el usuario la edita directamente, sin necesidad de volver
    a ejecutar el DSL.  Un backend alternativo podr\'{\i}a reproducir esta
    propiedad en su propio formato de salida.

\end{enumerate}

% ---------------------------------------------------------------------------
\chapter*{Recomendaciones}
\addcontentsline{toc}{chapter}{Recomendaciones}
\markboth{Recomendaciones}{Recomendaciones}

A partir del trabajo realizado se identifican las siguientes l\'{\i}neas de
extensi\'on para investigaciones futuras.

\begin{enumerate}

  \item \textbf{Soporte para m\'as funciones de celda.}  El DSL actual cubre
    \texttt{COUNTIF}, \texttt{COUNTA}, \texttt{SUM}, \texttt{IF},
    \texttt{AVERAGEIF} y operaciones de tiempo.  Funciones como
    \texttt{INDEX/MATCH}, \texttt{XLOOKUP}, \texttt{PIVOTBY} y las nuevas
    funciones de matrices din\'amicas de Excel~365 podr\'{\i}an a\~nadirse
    siguiendo exactamente el mismo patr\'on de extensi\'on descrito en la
    secci\'on~\ref{sec:extensibilidad}.

  \item \textbf{Backends alternativos.}  La independencia del backend respecto
    al AST hace posible implementar generadores de c\'odigo para otros
    formatos: ODS (LibreOffice Calc), CSV anotado, HTML con tablas,
    o incluso bases de datos SQL.  La capa del DSL no requiere ning\'un cambio.

  \item \textbf{Interfaz de usuario gr\'afica.}  El AST es una estructura de
    datos de primera clase en Lisp.  Una interfaz gr\'afica podr\'{\i}a visualizar
    y editar el AST directamente, permitiendo a usuarios no programadores
    configurar las reglas de coloreado y los c\'alculos sin escribir c\'odigo.

  \item \textbf{Verif\'{\i}caci\'on de tipos del DSL.}  El sistema actual no
    comprueba en tiempo de expansi\'on de macros que los s\'{\i}mbolos de columna
    referenciados en \dsl{exists} y \dsl{trange} existan en las tablas
    correspondientes.  Un verificador de tipos a nivel del AST detectar\'{\i}a
    estos errores antes de la generaci\'on del c\'odigo Python.

  \item \textbf{Generaci\'on incremental.}  El sistema genera el programa
    Python completo en cada ejecuci\'on.  Un modo incremental que detecte
    qu\'e tablas cambiaron y s\'olo regenere las hojas afectadas reducir\'{\i}a
    el tiempo de ciclo en workbooks con muchas hojas.

  \item \textbf{Plantillas reutilizables.}  Las plantillas de tabla definidas
    con \dsl{def-table} podr\'{\i}an exportarse a una biblioteca de plantillas
    est\'andar para los documentos m\'as comunes de la facultad (horarios de
    docencia, defensas, actas de ex\'amen), reduciendo el esfuerzo de
    configuraci\'on inicial para cada nuevo caso.

\end{enumerate}


% ---------------------------------------------------------------------------
%  REFERENCIAS
% ---------------------------------------------------------------------------
\bibliographystyle{unsrt}
\begin{thebibliography}{99}

\bibitem{fowler2010}
M. Fowler.
\newblock \textit{Domain-Specific Languages}.
\newblock Addison-Wesley Professional, 2010.

\bibitem{aho2006}
A.~V. Aho, M.~S. Lam, R. Sethi y J.~D. Ullman.
\newblock \textit{Compilers: Principles, Techniques, and Tools}.
\newblock Pearson/Addison-Wesley, 2\textsuperscript{a} ed., 2006.

\bibitem{steele1990}
G.~L. Steele.
\newblock \textit{Common Lisp: The Language}.
\newblock Digital Press, 2\textsuperscript{a} ed., 1990.

\bibitem{graham1993}
P. Graham.
\newblock \textit{On Lisp: Advanced Techniques for Common Lisp}.
\newblock Prentice-Hall, 1993.

\bibitem{seibel2005}
P. Seibel.
\newblock \textit{Practical Common Lisp}.
\newblock Apress, 2005.

\bibitem{keene1989}
S.~E. Keene.
\newblock \textit{Object-Oriented Programming in Common Lisp: A Programmer's
  Guide to CLOS}.
\newblock Addison-Wesley, 1989.

\bibitem{openpyxl}
E. Gazoni y C. Clark.
\newblock openpyxl --- A Python library to read/write Excel 2010 xlsx/xlsm/xltx/xltm files.
\newblock \url{https://openpyxl.readthedocs.io}, 2010--2024.

\bibitem{ooxml}
{ISO/IEC}.
\newblock \textit{ISO/IEC 29500: Information technology --- Office Open XML File Formats}.
\newblock International Organization for Standardization, 2008.

\bibitem{van-deursen2000}
A. van Deursen, P. Klint y J. Visser.
\newblock Domain-specific languages: an annotated bibliography.
\newblock \textit{ACM SIGPLAN Notices}, 35(6):26--36, 2000.

\bibitem{mernik2005}
M. Mernik, J. Heering y A.~M. Sloane.
\newblock When and how to develop domain-specific languages.
\newblock \textit{ACM Computing Surveys}, 37(4):316--344, 2005.

\bibitem{mccarthy1960}
J. McCarthy.
\newblock Recursive functions of symbolic expressions and their computation by
  machine, {P}art {I}.
\newblock \textit{Communications of the ACM}, 3(4):184--195, 1960.

\bibitem{clos1988}
D.~G. Bobrow, L.~G. DeMichiel, R.~P. Gabriel, S.~E. Keene, G. Kiczales y D.~A. Moon.
\newblock Common Lisp Object System specification.
\newblock \textit{ACM SIGPLAN Notices}, 23(S):1--142, 1988.

\bibitem{kiczales1991}
G. Kiczales, J. des Rivières y D.~G. Bobrow.
\newblock \textit{The Art of the Metaobject Protocol}.
\newblock MIT Press, 1991.

\bibitem{microsoft-excel}
{Microsoft Corporation}.
\newblock \textit{Excel for Microsoft 365 --- Reference Guide}.
\newblock Microsoft Docs, 2023.

\bibitem{ldmag2025}
R.~A. C\'apiro Colomar.
\newblock \textit{{LDMAG}: Lenguaje para el dise\~no de microaplicaciones de
  gesti\'on}.
\newblock Trabajo de Diploma, Facultad de Matem\'atica y Computaci\'on,
  Universidad de La Habana, 2025.

\bibitem{watt2004}
D.~A. Watt.
\newblock \textit{Programming Language Design Concepts}.
\newblock Wiley, 2004.

\bibitem{liskov1974}
B. Liskov y S. Zilles.
\newblock Programming with abstract data types.
\newblock \textit{ACM SIGPLAN Notices}, 9(4):50--59, 1974.

\bibitem{scott2015}
M.~L. Scott.
\newblock \textit{Programming Language Pragmatics}.
\newblock Morgan Kaufmann, 4\textsuperscript{a} ed., 2015.

\end{thebibliography}

\end{document}
